%!LW recipe=latexmk
%\documentclass[12pt,a4paper]{article}
\documentclass[12pt]{article}
\usepackage{amsmath}
\usepackage{amsfonts}
\usepackage{amssymb}
\usepackage[margin=1in]{geometry} % set margins
\usepackage{setspace} % for line spacing in document
\usepackage{float} % advanced positioning of figures and tables
\usepackage{graphicx} % including graphics
%\usepackage{hyperref} % allows for the insertion of urls
\usepackage[hidelinks]{hyperref}
\usepackage{bookmark} % creates PDF bookmarks
\usepackage{subfigure} % allows for two figures to be included in one figure
\usepackage{natbib} % customizes citations in the text
%\usepackage{biblatex} % customizes citations in the text
\usepackage{epstopdf} % ensures that eps files are converted to pdf so windows doesn't freak
\usepackage{mdwlist}
\usepackage{pdflscape}
\usepackage{epstopdf} % converts eps to pdf
\usepackage[section]{placeins} % keep figures within section
\usepackage{caption} %add captions below figures (use with below)
\usepackage{bbm}
\usepackage{eurosym}
\usepackage{comment}
\usepackage[title]{appendix}
%\usepackage{graphicx,showframe}
%\usepackage{subfigure}
%\usepackage{adjustbox}
\usepackage[mathscr]{euscript}
\usepackage{multirow} % multirow tables
\usepackage{rotating} % To display tables in landscape
\newcommand\fnote[1]{\captionsetup{font=footnotesize}\caption*{#1}} %custom captions
\usepackage{threeparttable, booktabs} % nice tables
\usepackage{tikz}
\usepackage{soul} %\hl
\usepackage{tabularx}
\usepackage{dutchcal}

\graphicspath{{/Users/drewvankuiken/Desktop/research/interconnection/drafts/proposal/pics}}

% https://tex.stackexchange.com/questions/375635/defining-column-width-using-pwidth-in-tabular-environment
% centered column with fixed width
\newcolumntype{C}{>{\centering\arraybackslash}p{1.3cm}}


%% OPERATORS
\DeclareMathOperator*{\argmax}{argmax}
% q^*
\newcommand{\qsi}{\ensuremath{q^*_{1}(\mu;\theta,\mathcal{P}_k)}}
\newcommand{\qst}{\ensuremath{q^*_{2,i}(\mu;\theta,\mathcal{P}_k)}}
\newcommand{\qstv}{\ensuremath{q^*_{2,t}(\mu;\theta,\mathcal{P}_k)}}
\newcommand{\qsj}{\ensuremath{q^*_{j}(\mu;\theta,\mathcal{P}_k)}}
\newcommand{\qs}{\ensuremath{q^*(\mu;\theta,\mathcal{P}_k)}}

% w(q*)
\newcommand{\wsi}{\ensuremath{w_{1}(\qsi)}}
\newcommand{\wst}{\ensuremath{w_{2}(\qst)}}
\newcommand{\wstv}{\ensuremath{w_{2}(\qstv)}}
\newcommand{\wsj}{\ensuremath{w_j(\qsj)}}
\newcommand{\ws}{\ensuremath{w(\qsi,\qst)}}

% O(q^*)
\newcommand{\os}{\ensuremath{\mathcal{O}(\qsi,\qst;\mathcal{P}_k)}}

% W^* and O^*
\newcommand{\ew}{\ensuremath{w^*(\theta, \mathcal{P}_k)}}
\newcommand{\ewm}{\ensuremath{w^*(\theta, \mathcal{P}_m)}}
\newcommand{\eo}{\ensuremath{\mathcal{O}^*(\theta, \mathcal{P}_k)}}
\newcommand{\eom}{\ensuremath{\mathcal{O}^*(\theta, \mathcal{P}_m)}}

% Lagrangian
\newcommand{\Lagr}{\mathcal{L}}

% removes the padding that is normally placed between sections
%\usepackage[compact]{titlesec}
%\titleformat{\section}{\large\bfseries}{\thesection}{0.5em}{}
%\titleformat{\subsection}{\normalsize\bfseries}{\thesubsection}{0.5em}{}
%\titleformat{\subsubsection}{\small\bfseries}{\thesubsubsection}{0.5em}{}
%\titleformat{\subsubsubsection}{\small\bfseries}{\thesubsubsubsection}{0.5em}{}
%\titlespacing{\section}{0pt}{0.25in}{*0}
%\titlespacing{\subsection}{0pt}{0.25in}{*0}
%\titlespacing{\subsubsection}{0pt}{0.25in}{*0}
%\titlespacing{\subsubsubsection}{0pt}{0.25in}{*0}

\def\sectionautorefname{Section}

%\usepackage[backend=biber]{biblatex}
%\hypersetup{citecolor=black}
%\bibliography{../presentations/bib_interconnect.bib}
%\bibliographystyle{aea.bst}

\title{Infrastructure as Market Design: Location and Technology Choice on the Electric Grid\footnote{I am grateful to Jon Williams, Andy Yates, Gary Biglaiser, Marco Duarte, and Valentin Verdier for their help. I thank seminar and conference participants for helpful comments on this paper. Department of Economics, University of North Carolina at Chapel Hill, \href{mailto:drew.van.kuiken@unc.edu}{drew.van.kuiken@unc.edu}.}}

\author{Drew Van Kuiken}

\date{\today}
%\date{December 12, 2024}


\begin{document}

\maketitle

\begin{abstract}

%\noindent When a new renewable generator enters the grid, it takes up scarce transmission capacity and often displaces costly natural gas. Using prices to internalize entrant externalities is difficult in this setting: congestion is unpredictable and its costs are spread widely; the emissions savings due to a new renewable generator change by the minute; and new transmission capacity has benefits for everyone on the grid. Grid operators use two policy instruments to manage entry to the grid: location-specific entry fees and grid upgrade decisions. Different grids manage entry very differently -- Texas, the setting for this paper, imposes minimal entry fees and lets generators manage their congestion exposure themselves. The Mid-Atlantic levies very high entry fees. Both grids use myopic engineering tests to determine when grid improvements should be made. I study how grid operator entry policies shift the level and location of entry. My setting is Texas, where I build a simplified model of the grid using new and detailed data on generator locations, cost curves, and grid capacity. To answer this question, I estimate a new dynamic, spatial, and structural model of generator entry from 2018 to 2025, a period when dramatic shifts in renewable capital costs and tax treatments led to increased entry for renewables. I find XXX. In counterfactuals, I assess how entry would shift if Texas were to implement the entry policies that the Mid-Atlantic uses and if Texas were to use forward-looking grid improvement policies. I find that YYYY.
\onehalfspacing 
\noindent If benefits to infrastructure vary by firm, shifting infrastructure availability will change the equilibrium level and location of entry. I study this channel in the context of the Texas electric grid, where electric transmission is scarce and renewable resource quality varies by region and technology. I solve a dynamic entry game for renewable generators, capturing spatial and technological entry margins. Embedded within my entry game is a physically realistic model of electricity dispatch, where electricity flows and prices are endogenous. The model is based on a new dataset (2018-2024) with generation, transmission capacity, and demand information. I find that limited transmission capacity and spatially concentrated wind resources suppress wind entry, while solar generators respond to limited transmission capacity by entering less sunny regions with fewer transmission constraints. I explore two counterfactuals -- one where entrants pay for grid upgrades, and one where grid operators run a market for new transmission -- under fixed demand and data-center-driven demand growth. I find that the market for new transmission generates more welfare than the Texas or entrant-funded upgrade policies, and that both non-market policies yield inefficiently low and spatially distorted entry choices.

%I implement a commonly-used transmission policy where entrants must pay for grid upgrades to facilitate their own entry and a policy where grid operators set up a market in which entrants, incumbents, and consumers reveal their willingness to pay for new transmission. I find the market lowers congestion, lowers prices, increases entry in high resource quality regions, and constructs more transmission than Texas' baseline policy or the upgrade-on-entry policy. Finally, I consider how infrastructure policies deal with data center driven demand growth, adding 1 GW of demand in different regions on the grid. I find that the market for transmission greatly outperforms existing policies. 


%Counterfactuals consider alternative transmission financing schemes: entrant funding of new transmission, which significantly limits the amount of generation on the grid; and a forward market for new transmission, which is welfare improving relative to the Texas and entrant-funded systems. Finally, I test how each financing scheme responds to the addition of a 1 GW data center in a random location, showing that a forward market vastly outperforms existing policies. These results illustrate how product positioning can be endogenous to infrastructure provision: transmission shapes both firms' equilibrium locations and the technology mix that emerges.

%What is the best way to provide new transmission when generator entry is endogenous? To answer this question, I assemble a new dataset on generation, transmission capacity, and demand for the Texas electric grid. I use these data to construct a physically realistic model of electricity dispatch, where electricity flows and prices are endogenous. I then solve a dynamic entry game for renewable generators, capturing spatial (where wind and solar generators choose to enter the grid) and technological (the amount of wind and solar on the grid) margins. I find that limited transmission capacity and spatially concentrated wind resources suppress wind entry, while solar generators respond to limited transmission capacity by entering less sunny regions with fewer transmission constraints. Counterfactuals consider alternative transmission financing schemes: entrant funding of new transmission, which significantly limits the amount of generation on the grid; and a forward market for new transmission, which is welfare improving relative to the Texas and entrant-funded systems. Finally, I test how each financing scheme responds to the addition of a 1 GW data center in a random location, showing that a forward market vastly outperforms existing policies. These results illustrate how product positioning can be endogenous to infrastructure provision: transmission shapes both firms' equilibrium locations and the technology mix that emerges.

% mention market structure/infra non neutrality



\vspace{0.25in}

%\noindent\textbf{Keywords}: Renewable

\vspace{0.1in}

%\noindent\textbf{JEL Codes}: L1, L93
\end{abstract}

\break
	
\onehalfspacing
%\doublespacing


Economists have long recognized that infrastructure is important for welfare, but figuring out the marginal benefit of infrastructure is difficult \citep{Fernald1999,DurantonTurner2011,BrancaccioKalouptsidiPapageorgiou2024}. Effects are highly nonlinear, both in levels and locations. Prior literature suggests a link between infrastructure and market structure: when infrastructure availability changes, some firms benefit more than others \citep{SniderWilliams2015,Trew2020}. This paper begins by observing that if benefits to infrastructure vary by firm, shifting infrastructure availability will change the equilibrium level and location of entry. Thus welfare changes depend on how infrastructure shifts the composition of firm entry, particularly when firms differ in their productivity, costs, or consumers they serve. 

I study how infrastructure policy shifts equilibrium level, location, and technological composition of entry in the context of the Texas electric grid, where these dynamics are particularly relevant. Texas has ample wind and sunlight resources. It also builds new electric transmission infrequently. These two features of the market came to a head during the 2018 to 2024 period, when fixed costs for wind and solar generators dropped and both technologies entered rapidly. Grid capacity remained fixed throughout the period and grid congestion increased. The incidence of this congestion was unequal: generators stuck behind congested transmission lines saw revenues fall, while generators in uncongested regions of the grid earned steady revenues. The geography of resource quality made this particularly punishing for wind generators -- while solar generators had access to reasonably high quality sunlight throughout the state, high quality wind resources were located in the panhandle and central Texas, gated behind congested transmission elements. This makes Texas an ideal case to test how infrastructure availability shifts the level and location of entry when products are differentiated.

To test the effect of infrastructure on entry, I solve a dynamic entry game among renewable generators, where generators differ in their technology and their location on the grid. Entrants are forward-looking and make decisions based on the expected evolution of entry both in their own region and in other regions. I find that limited transmission capacity interacts with wind's lack of geographic flexibility to reduce wind entry, while solar generators sacrifice sunlight quality to pursue available transmission capacity. 

I then ask whether Texas' transmission policy leads to an efficient allocation of generation on the grid. I begin by implementing a transmission policy used throughout the rest of the U.S., where entrants must pay for grid upgrades to facilitate their own entry. I find that this policy imposes inefficiently high costs on entrants, leading to cascading withdrawal and lower welfare. I then test the effect of setting up a market where entrants, incumbents, and consumers reveal their willingness to pay for new transmission. I find this policy lowers congestion, lowers prices, increases entry in high resource quality regions, and constructs more transmission than Texas' baseline policy, resulting in a XXX\% increase in welfare. Finally, I consider how infrastructure policies deal with data-center-driven demand growth, adding 1 GW of demand in different regions on the grid. I find that the market for transmission greatly outperforms existing policies in this case, increasing welfare by YYY\% over the Texas baseline. 

Together, my results suggest three takeaways: (i) infrastructure is non-neutral towards differentiated technologies, shaping firms' equilibrium locations and technology choices on the electric grid; (ii) without a market for new transmission, realized entry choices are inefficiently low and spatially distorted; (iii) coming data center demand is likely to amplify existing entry inefficiencies, but better infrastructure policy can leverage increased demand to build more transmission and lower electricity prices.

The analysis is built on a newly assembled dataset, featuring detailed information on generation, transmission capacity, and demand for the Texas electric grid. Some transmission capacity is unobserved. I develop a procedure to estimate unobserved transmission capacity based on transmission voltage and county-level generation and demand. I pair this information with yearly data on renewable fixed costs, including capital costs and tax rates, and resource quality, which measures the wind and sunlight available for generators to turn into electricity at the county level. As I describe below, my model is able to capture how resource quality variation shifts entry decisions differentially by technology.

Entry depends on profits and profits depend on prices and quantities sold. I construct a physically realistic model of electricity dispatch, where electricity flows, prices, and quantities sold are endogenous. The model captures the physical realities that generators must contend with: when transmission is congested, generators in counties upstream of the congestion cannot send their electricity to the highest demand regions. Prices in these counties fall as a result, often reaching levels near zero.\footnote{This occurs because renewable marginal costs are zero and electricity prices reflect the marginal cost of the marginal generator in a given county.} Prices in counties downstream of the congestion rise.\footnote{If no transmission is congested, prices are equal across all counties.}

The entry model has 40 regions, which span the service territory of ERCOT, the grid operator for Texas. Each period, a fixed number of potential entrants arrive in each region and decide on entry. Regions vary in the quality of resource they have available and the expected price that entrants will earn from producing in that region, among other factors. Resource quality is exogenous and fixed; regional prices and quantities are endogenous. Entrants make decisions based on expected profits, weighed against the fixed costs of entry for a given region and a cost shock.

Prices on the electric grid depend on the flows of electricity throughout the network. Solving the entry game while accounting for evolving wind and solar generation in all 40 regions would require a computationally infeasible state space. Instead of solving the full game, I solve a simplified game, where entrants form beliefs about other-region entry based on equilibrium entry policies and integrate over potential spatial allocations of generation that cost shocks could justify. This preserves the network structure of prices: entry in one region affects prices everywhere. In equilibrium, entrant beliefs align with the stationary distribution of generation allocations, meaning that entrant expectations about prices are accurate within the model. 

The entry model confirms that wind generators are more heavily affected by congestion than solar generators: in a partial equilibrium exercise, I find that setting congestion costs equal to zero increases wind entry by 8.9\%, primarily in the regions where wind generation already exists, while solar entry increases by less than 2\%.\footnote{The spatial allocation of solar instead shifts, as generators enter in higher quality sunlight regions.} Congestion incidence falling unequally on technologies does not mean that Texas' policy is inefficient, however. 

Two potential inefficiencies emerge from the entry model: entry has network spillovers -- transmission is scarce and incumbents have no priority in using it; and transmission is underprovided, even as values for new transmission are heterogeneous across locations and technology types. Entry behavior in Texas suggests that alternative infrastructure policies could increase efficiency and welfare. To test whether this is true, I implement the upgrade-on-entry and market for transmission policies. I find that upgrade-on-entry roughly targets both margins of inefficiency, but the costs of entry are too high and borne exclusively by entrants when new transmission also benefits incumbents. Since upgrade costs are shared within entry cohorts, the consequences of adopting this policy are significant: entrants often withdraw instead of paying their upgrade cost allocation, leaving upgrade costs split among fewer entrants. In some cases, withdrawal decisions cascade. This reduces welfare by XXX\% relative to the Texas baseline. I find that the market for new transmission instead mitigates these distortions and improves upon the Texas baseline with respect to efficiency and welfare. 

In the near term, demand from data centers is poised to cause new challenges for the grid. Whereas electricity demand has been mostly flat over the past 25 years, Texas grid operators expect demand to grow by XXX\% over the next YYY years. In a second counterfactual, I examine how the Texas baseline, the upgrade-on-entry policy, and the market for new transmission all respond to the addition of a 1 GW data center to the grid. I find....

%I find that wind and solar generators trade off resource quality and transmission availability when deciding on entry. I find this tradeoff is considerably more pronounced for wind generators. In the lowest resource quality regions, solar generators face a \$0.02/MWh discount due to congestion (relative to the uncongested price), while in the highest quality regions, they face a \$0.24/MWh discount. For wind, the lowest quality regions face a \$0.46/MWh discount, while the highest quality regions face a \$1.19/MWh discount.\footnote{Wind generators face a steeper discount than solar generators in all circumstances because wind generators produce during times of the day when demand is lower, driving prices down and curtailment up. Curtailment rises with the low demand because wind production is geographically concentrated: when the wind blows, transmission lines cannot carry the production of the wind generators.} In a partial equilibrium exercise, I find that setting congestion costs equal to 0 increases wind entry by 8.8\%, primarily located in regions with existing wind resources. 

%More than partial equilibrium is available however: the dispatch model and entry model permit evaluating the welfare effects of different infrastructure policies. I investigate how two different policies compare to the Texas baseline: an ``invest-and-connect'' policy, which is similar to ones used throughout the rest of the U.S.; and a market for transmission policy, in which generators (entrants and incumbents) and consumers (truthfully) reveal their willingness-to-pay for new infrastructure and the grid operator selects the upgrade package that maximizes bids minus cost. %The market for transmission policy operates similarly to open season mechanisms from gas pipeline development (\cite{Kavulla2026}). 

\subsubsection*{Literature Review}

I contribute to a long-running literature in economics studying strategic entry and the positioning of firms and products (\cite{Mazzeo2002}; \cite{Seim2006}; \cite{Jia2008}; \cite{Ryan2012}; \cite{Sweeting2013}). This literature studies how firms choose the products they sell and the locations they serve. The electric grid is uniquely well suited to understand endogenous technology choice and siting decisions: wind and solar production profiles depend on exogenous environmental factors and locations are differentiated based on shared transmission infrastructure. I leverage this setting to estimate a dynamic entry model with endogenous congestion and transmission provision. The model captures how infrastructure policy changes the equilibrium assortment of wind and solar generation. I show that geographically flexible solar projects can move towards demand centers when infrastructure is scarce, while wind projects are constrained by the environment, resulting in under-entry. I am additionally able to use this model to consider how alternative grid management policies would shift the spatial allocation of firms, the equilibrium technology mix, and welfare. 

This paper adds to a growing literature concerned with the electric grid. Grid management is complicated and economists have only recently begun to model the processes that RTOs use to govern entry and competition on the grid (\cite{Wolak2015}; \cite{GowrisankaranReynoldsSamano2016}; \cite{Cicala2022}; \cite{Mays2023}; \cite{JohnstonLiuYang2024}; \cite{Boatwright2025}). I add to these papers the first embedding of a realistic dispatch model within a dynamic entry game. I use data from existing generators and the offers they submitted to the grid operator, a novel representation of the grid, and endogenize congestion in a way that is new to the literature. Previous work has used representations of the grid (\cite{ArkolakisWalsh2023}; \cite{Golden2026}), but focused on the grid as a series of bipartite connections. This paper instead models the physics that govern where power flows on the grid. This allows me to model how congestion on one grid segment affects prices throughout the grid. More broadly, it allows me to credibly estimate the dynamic incentives firms face when choosing where they should enter the grid and consider how different grid management policies would shift the level and location of renewable entry. 

The benefits and underprovision of transmission has been studied in a variety of settings, including the Midwest, Texas, Chile, and the continental U.S. (\cite{BorensteinBushnellStoft2000}; \cite{FellKafineNovan2021}; \cite{GonzalesItoReguant2023}; \cite{ArkolakisWalsh2023}; \cite{Hausman2024}; \cite{HollandMansurYates2024}; \cite{Boatwright2025}; \cite{Golden2026}). These papers show that transmission increases allocative efficiency, decreases emissions, and increases long-run investment in renewable capacity. Despite widespread agreement as to the value of new transmission, underinvestment in transmission persists. On this point, my paper is closest to \cite{Boatwright2025}, who also studies entry to the grid in Texas. She estimates a two-stage model of developer applications and project completion, showing how observed congestion exposure affects entry into the interconnection queue and eventual construction. In contrast, I model the dynamic feedback between entry, dispatch, congestion, and transmission planning: firms' siting decisions change future grid conditions, and grid upgrades respond to the resulting state of the network. By endogenizing the transmission planning process, I connect this literature to an older literature that theorized about transmission planning policies for RTOs (\cite{Leautier2000}; \cite{Vogelsang2001}; \cite{JoskowTirole2005}; \cite{HoganRosellonVogelsang2010}). These papers recognize the difficulty in incentivizing efficient expansion of the transmission grid, showing how underinvestment in transmission arises. I capture this underinvestment using my model of centralized transmission planning. My main contribution to these literatures is to estimate the costs of centralized transmission planning when entry is endogenous. I also propose a new mechanism of transmission funding that accounts for the strategic entry of renewable energy and which leads to larger investment in the grid.

The rest of the paper is structured as follows: Section \ref{sec:background} introduces the setting and data. Section \ref{sec:ps_and_qs} discusses the dispatch model, including transmission line capacity estimation and the commitment model. Section \ref{sec:entry_model} introduces the dynamic entry model, while Section \ref{sec:estimation} describes how I estimate the model. Section \ref{sec:results} describes the entry model results and Section \ref{sec:counterfactuals} introduces the counterfactuals and the results I find. Section \ref{sec:conclusion} concludes.

\section{Background and Data}
\label{sec:background}
The late 1990s brought a wave of electricity market restructuring to the grid. Whereas utilities had historically been in charge of building new electric generation, most grids in the U.S. decided to let private investors determine where and when new generation would be built.\footnote{Recent research suggests that power plant operations have become more efficient following the restructuring, while the effect on prices has been muted (\cite{BorensteinBushnell2015}).} Utilities remained control of the grid itself: where transmission lines go and how much electricity they can carry. This paper studies how private generation and public infrastructure interact to determine where generators enter on the grid.

At the time of the switch, ERCOT ran on dispatchable electricity. Figure \ref{fig:gen_map_2000} shows generators that operated in ERCOT in 2000. At the time, generators clustered around major urban areas and industrial load zones like the Gulf Coast. In a grid dominated by dispatchable resources, transmission primarily exists to smooth demand shocks: if a heat wave hits Dallas, generators in Houston can turn on and send electricity for air conditioning and vice versa. Outside of load zones, transmission plays a smaller role. Instead, counties have one or two generators that meet local demand. 

\begin{figure}[htbp]
    \centering

    \vspace{0.4em}
    \includegraphics[width=0.92\textwidth]{../figures/paper/texas_generation_shared_legend.pdf}

    \subfigure[2000\label{fig:gen_map_2000}]{%
        \includegraphics[width=0.47\textwidth]{../figures/paper/texas_generation_2000_eia/texas_generation_facilities_eia_2000.pdf}%
    }%
    \hfill%
    \subfigure[2025\label{fig:gen_map_2025}]{%
        \includegraphics[width=0.47\textwidth]{../figures/paper/texas_generation_2025_eia/texas_generation_facilities_eia_2025.pdf}%
    }

    \caption{ERCOT generating capacity, 2000 and 2025}
    \label{fig:gen_maps}
    \caption*{\footnotesize
        \textit{Notes:} Each point represents a generator. Marker area
        scales with nameplate capacity, and color denotes the facility's dominant
        fuel category. ``Other'' includes nuclear, hydroelectric, biomass,
        storage, and other fuels. Pale-blue shading denotes counties in ERCOT;
        the light-red line traces the resulting county-level ERCOT boundary.
        Data come from the historical EIA-860A/B inventories for 2000 and the
        EIA-860 2025 early release.
    }
\end{figure}

Figure \ref{fig:gen_map_2025} shows generators that operated in ERCOT in 2025. The rise of renewables in ERCOT has been remarkable. Whereas coal and gas represented 92\% of capacity in 2000, they represent 48\% in 2025. Renewables represent 40\%.\footnote{Renewables produced 37\% of energy consumed in ERCOT in 2025 because they only infrequently produce at full capacity.} The majority of renewable additions have come since 2017, following dramatic decreases in solar fixed costs and generous tax incentives for the development of wind power. Texas' renewable resources are located primarily in the western and northwestern part of the state; good solar resources are available in the far south as well. The highest-quality renewable resources are located far from ERCOT's demand centers. 

The shift towards renewables has created issues for the existing transmission paradigm in ERCOT. Because of renewables, transmission today amplifies the variance of supply.  Grid operators are left to decide on what the proper amount of transmission is. Building infinite transmission would allow the grid to capture low price electricity produced by renewables, but the lines would sit empty much of the time. Ignoring the opportunity would leave Texas' considerable solar and wind resources untapped, with prices and emissions higher than they need to be. 

Existing policy towards transmission varies. In ERCOT, operators allow renewables to connect freely to the grid, whether or not the transmission capacity exists for them. Renewable generators take on the risk of curtailment as a consequence, which is when a generator has to the throw away the electricity it produces because it cannot reach demand. In PJM, which I return to in counterfactuals, grid operators ask generators to upgrade the grid themselves such that they will be able to transmit the electricity they produce to load zones in the vast majority of situations. 

ERCOT's choice -- letting renewables interconnect freely and curtailing extra production -- has come with consequences. Congestion on the grid arises when supply (largely from renewables) outstrips available transmission capacity. The entry of renewables between 2018 and 2024 led to unprecedented congestion, with yearly congestion costs reaching a high of \$2.8B in 2022 (\cite{PotomacEconomics2023}). Transmission upgrades have not kept pace with the renewable capacity and transmission availability has decreased as a result.

Where and when renewables enter depends on transmission availability. With infinite transmission, generators would be free to locate exclusively based on resource quality, finding the sunniest and windiest spots possible. As transmission availability shrinks, generators must trade off between resource quality and deliverability. Further, while transmission availability is common among solar and wind generators, resource quality is not. Texas' wind resources are localized: utility-scale wind is most feasible in the Panhandle, central Texas, and the far south. The solar resource quality gradient is more gentle. While the best solar resources are located in the south and far west parts of the state, solar resource quality is fairly good throughout the state, including nearby to load zones in the east.

Figure \ref{fig:resource_headroom_entry} shows how solar and wind generators navigate the resource quality-transmission capacity tradeoff in ERCOT. The heatmap shows entry rates at the county-technology-year level. Observations are broken into quartiles based on resource quality and transmission capacity. The solar figure shows that 15\% of entrants in the high resource quality, second quartile transmission capacity bin entered between 2018 and 2024. Similar entry rates persist for third quartile-third quartile cells and high capacity, low resource quality cells. The tradeoff is less stark for wind generators, for which low resource quality cells are near-impossible to profitably operate in. Instead, wind generators cluster around high resource quality locations, even when transmission capacity for those locations is limited. Overall, wind generators enter in 5\% of observations, while solar generators enter in 8\%.

\begin{figure}[htbp]
    \centering
    \includegraphics[width=\textwidth]{../figures/entry_patterns/pdf/resource_headroom_entry_rate_surface.pdf}
    \caption{Renewable entry across resource quality and inherited transmission capacity}
    \label{fig:resource_headroom_entry}
    \caption*{\footnotesize
        \textit{Notes:} Cells report the share of county-technology-year
        observations with at least one entrant, pooled over 2018--2025.
        Resource quality and predetermined physical capacity are divided into
        quartiles within each technology-year. Solar and wind use a common color
        scale. Blank cells contain fewer than ten county-year observations.
    }
\end{figure}


% Figure \ref{fig:centroids} shows equilibrium entry by technology and year in ERCOT. When capacity was abundant, solar was able to seek out better resource quality in the west. As congestion increased, solar moved eastward, locating closer to demand. The resource quality gradient for wind resources is much steeper. Good wind resources are primarily located in the Texas panhandle and central Texas. As congestion rose in ERCOT, wind's equilibrium entry location did not shift to the same extent as solar's did. 


\subsection{Policy}
\label{sec:policy}

ERCOT's policy with respect to entry -- letting generators freely join the grid with active curtailment -- is uncommon in the United States.\footnote{Generators must still cover point of interconnection fees, see XXX.} Most other grids like PJM use an ``invest-and-connect'' process. When a generator enters the PJM grid, PJM initiates a study process to identify any overloads that the generator might cause and suggests a series of improvements to the grid that the generator should make to avoid grid congestion in summer peak, winter peak, and light demand scenarios. The generator must then make investments to improve the grid before it can enter. The scope of the upgrades that a generator must make thus depends on the location it chooses: locations with poorer local transmission capacity are more likely to trigger expensive upgrades. The grid is public infrastructure and investments are non-rival up to grid capacity.

The invest-and-connect policy is costly for renewable plants. In PJM, generators wait for 60 months on average between the time they enter the study process and when they successfully connect (\cite{RandEtAl2025QueuedUp}). Renewables typically face higher entry fees than natural gas as well; in 2023, solar entry fees averaged \$253/kW, wind entry fees averaged \$136/kW, and natural gas entry fees averaged \$24/kW (\cite{SeelEtAl2023PJMInterconnectionCosts}). In principle, these costs force generators to internalize the value of scarce transmission capacity, a signal that is missing from ERCOT's market design. 

Once on the grid, generators face less congestion in PJM than those in ERCOT do. Table \ref{tab:congestion-costs} reports yearly congestion costs in each grid. Congestion costs in PJM were lower than those in ERCOT for 2019 through 2024 despite PJM being a much larger grid than ERCOT (\cite{ShreveEtAl2025TransmissionCongestion}). 

\begin{table}[htbp]
\centering
\caption{Transmission Congestion Costs and Transmission Miles}
\label{tab:congestion-costs}
\begin{tabular}{lrrrrrrrr}
\toprule
 & \multicolumn{7}{c}{Congestion costs (\$M)} &  \\
\cmidrule(lr){2-8}
RTO/ISO & 2018 & 2019 & 2020 & 2021 & 2022 & 2023 & 2024 & Transmission miles \\
\midrule
ERCOT & 1{,}260 & 1{,}260 & 1{,}400 & 2{,}100 & 2{,}800 & 2{,}400 & 1{,}900 & 55{,}000 \\
PJM   & 1{,}310 & 583     & 529     & 995     & 2{,}500 & 1{,}068 & 1{,}754 & 88{,}333 \\
\bottomrule
\end{tabular}

\vspace{0.5em}
\begin{minipage}{0.95\textwidth}
\footnotesize \emph{Notes:} Congestion costs are in millions of dollars and come from Table B1 of \citet{ShreveEtAl2025TransmissionCongestion}. Transmission mileage comes from PJM and ERCOT fact sheets.
\end{minipage}
\end{table}

Grid upgrades are centrally planned across the U.S.\footnote{Though the invest-and-connect regime also improves the grid when generators enter.} Generally, utilities are responsible for identifying and proposing new upgrades, while grid operators retain final approval authority. Utilities earn a regulated return on these grid investments, with rates subject to review by Public Utility Commissions. Recent studies have found that transmission planning is falling far behind generation growth and, more recently, load growth. Today, over 90\% of upgrades are focused on reliability (i.e., avoiding transmission line overloads), instead of economic benefits (i.e., reducing congestion costs). 

Funding for grid upgrades generally comes from consumers. ERCOT allocates costs based on a beneficiaries pay system, where the end-users who see lower prices as a result of new transmission fund the transmission upgrades. PJM mostly follows a beneficiaries pay system, though generators are responsible for grid upgrades that are identified in the interconnection process. Upgrades are triggered by a Production Cost Savings Test, which calculates whether a proposed investment would have paid for itself in saved congestion costs over the past year. PJM uses a similar test that blends the Production Cost Savings Test with a more regionally targeted savings threshold.\footnote{I need to pin down exactly how PJM chooses between assigning costs to generators versus using this centralized grid upgrade process. It's a little confusing. I think I can implement either method in counterfactuals, and neither method is going to replicate first-best transmission allocation.}

%With data center additions further stressing the transmission grid, policymakers have started to consider new options for growing the grid. Of particular interest is adapting the ``Open Season'' framework from gas pipeline development (\cite{Kavulla2026}). With an open season policy, a grid operator would develop a plan to increase capacity on the transmission grid. It would then solicit bids from demand, generators, and data centers that represent a willingness to pay for that increase in capacity. If the sum of the bids outweights the cost of the improvement, the improvement would be built.\footnote{Need to figure out if I'm proposing combinatorial auctions or not.} This idea does a much better job recovering the planner's solution to generator siting and transmission construction and I return to it in the counterfactuals.

\subsection{Grid}
\label{sec:grid}

Every 5 minutes, grid operators solve a constrained optimization problem: given demand for electricity, find the least-cost dispatch of electricity generation subject to transmission capacity. The operator solves this problem for each node across the entire grid. Nodes, typically represented by physical substations, are the point at which electricity moves from generators to the transmission grid and the point at which demand is settled for a small area. There are several thousand nodes in ERCOT. In this paper, I will model the grid operator's problem at the county level, aggregating all of the nodes that are within that county into one decision point. I will also model the problem at the hourly level to match the resolution of the demand data that is available to me. 

The grid operator is responsible to telling which generators to produce electricity each period. To facilitate this process, generators report their marginal cost of production to the grid operator.\footnote{I assume generators report their marginal costs honestly. Independent market monitors are appointed to each RTO to ensure that reports are honest. Market monitors are able to access internal generator data and grid operator data to ensure that generators are not manipulating their bids.} The grid operator collects generator bids and finds the point at which supply meets demand assuming no transmission constraints bind. This sets the system cost of energy, which every generator that produces electricity would earn for its production absent congestion. 

\begin{figure}[htbp]
    \centering
    \includegraphics[width=0.88\textwidth]
        {../figures/paper/grid_dispatch_illustration/nueces_offer_stack_20250514_1100.pdf}

    \caption{Supply, commitment, and congestion in Nueces County}
    \label{fig:nueces_offer_stack}

    \caption*{\footnotesize
        \textit{Notes:} The figure shows ERCOT SCED energy-offer blocks for generating resources mapped to Nueces County at 11:00:15 a.m. on May 14, 2025. Color denotes fuel; solid blocks were dispatched and pale blocks were not dispatched. Dotted gas blocks represent the $[0,\mathrm{LSL}]$ minimum-output requirements of committed generators, shown at $-\$50$/MWh to represent fixed injections. Offers above LSL are the generators' reported SCED offer curves. The solid red line is the SANTACRU\_ALL settlement-point price for the corresponding 15-minute interval; the dashed line is ERCOT system lambda for the displayed SCED interval. Committed thermal minimum output was 525 MW, while only 4 MW of thermal generation above LSL was dispatched. Wind output was 164 MW against an HSL of 303 MW. Data come from ERCOT's 60-day SCED disclosures and real-time
        settlement-price reports.
    }
\end{figure}

Balancing supply and demand on the grid is a high-stakes endeavor. If too much or too little supply comes online, stress on grid equipment increases and blackouts become more likely. The difficulty of this task is multiplied by intermittent renewable production: when clouds roll over a solar farm, the grid operator needs to adjust generation elsewhere on the grid to maintain the balance of supply and demand. To combat this problem, generators can enter a day-ahead market to commit themselves to produce 24 hours after they bid. Alternatively, grid operators in ERCOT can force dispatchable resources like coal and natural gas to produce 24 hours in advance. This process is called commitment.\footnote{PJM, along with most other grids in U.S., use capacity markets to ensure sufficient dispatchable generation exists on the grid each period. The economics of capacity markets lie outside the scope of my article because they play only a limited role in renewable entry decisions.} The volume of resources committed each period depends on forecasted demand, forecasted renewable supply, and the net generation needs at each node on the grid. Committed resources produce before any renewable production occurs, shifting the nodal supply curve to the right. 

Constraints are based on the transmission grid. Each link in the transmission grid has a total amount of electricity that it is able to carry. Counties in ERCOT are generally connected by high-voltage alternating current (HVAC) lines that move electricity from the resource-rich, low-demand regions in the western part of the state to high-demand population centers in the east. Congestion occurs once least-cost dispatch of electricity requires moving more electricity over a transmission line than that line's capacity can manage. 

Once lines become congested, the grid operator is forced to dispatch electricity from nodes with higher marginal cost than the system cost and curtail electricity production from nodes with lower marginal cost than the system cost. In nodes where generators are curtailed, the local price (LMP) falls below the system price of energy. LMPs indicate what the marginal generator in a node is willing to produce electricity for. In nodes on the demand side of congestion, prices will rise above the system price of energy, as marginal generation in that node has to turn on to meet demand. 

Figure \ref{fig:nueces_offer_stack} shows an example of the dispatch process in action. The figure shows the supply curve for Nueces county, TX, at 11 A.M. on May 14, 2025. Nueces is in southern Texas and has natural gas, wind, and storage generators. The system price of electricity during this period was \$25.31, while the LMP in Nueces was -\$21.33. The grid operator had committed three natural gas generators to produce 529 MW of electricity during this period. Those units can be seen on the left side of the supply curve.\footnote{Committed generators earn a price settled in day-ahead markets and often increased by ERCOT's reliability unit commitment process, which is intended to compensate generators for providing commitment instead of bidding into real-time electricity markets.} Export constraints meant that wind got curtailed relatively quickly in Nueces. Only 164MW of 303MW total available generation was produced. The remaining wind generation was curtailed. As a result, wind was the marginal generator in Nueces and the LMP dropped to wind's marginal offer.\footnote{Wind marginal prices are often negative due to tax incentives and power purchase agreements.}

\subsection{Data}
\label{sec:data}

An unusual feature of the electric grid is that, once interconnected, generators' scope for strategic decision-making is limited. Prices are set based on marginal cost, which is monitored for honest bidding, and capital improvements are rare.\footnote{Of 634 renewable generators operating in ERCOT between 2018 and 2024, only 11 (2\%) reported nameplate capacity increases greater than 10MW.} I abstract away from capital improvements in this paper. Thus, understanding entry decisions requires understanding the price that generators will earn in the market, the quantity of electricity that generators will be able to sell in the market, and fixed entry costs.

I compile three unique datasets to study entry empirically: an hourly production panel, which contains information on generators, their supply curve, their location, and demand; a yearly infrastructure panel, which contains information on transmission lines, substations, and upgrades; and a yearly fixed cost panel, which contains information on capital costs, taxes, and other fixed cost shifters. Transmission capacity is unobserved; in Section \ref{sec:network}, I describe the process I use to recover transmission capacity and validation checks I perform to ensure my capacities are reasonable. I describe the data construction for this project further in Appendix XXX. 

\subsubsection{Production Panel}
\label{sec:production}
ERCOT is unique in the level of production data that it makes available.  Using Seasonal Assessment of Resource Adequacy (SARA) and Monthly Outlook for Resource Adequacy (MORA) reports, I compile a monthly panel of generators that includes information on generator county, technology, interconnection date, and nameplate capacity. I link this panel to ERCOT Security Constrained Economic Dispatch (SCED) data at the generator level. SCED data contains the supply curves that each generator submits to the grid operator for constrained optimization, along with the amount of electricity they sold in a period. The economics literature typically treats renewable marginal costs as zero, which is borne out in the SCED data.\footnote{In fact, renewables often submit negative supply bids, since they earn tax credits for producing.} 

\input{../figures/paper/production_summary_by_weather_zone.tex}

County-level demand data for ERCOT is unavailable. I construct county-level demand based on ERCOT's hourly zonal load data (NP6-345-CD). The zonal load data is given for each of the eight ERCOT weather zones. I mapped counties in ERCOT to weather zones by hand. I combine this data with two products from the U.S. Energy Information Administration (EIA): data on electricity delivered to customers in 2020 (EIA Form 861 Part C, Delivery Companies) and utility service territories, which reports the counties served by each utility that delivered electricity in ERCOT in 2020 (Form 861, Service Territories).\footnote{County demand allocations are robust to using different years of EIA Form 861 data. Form 861 Part C begins in 2020. 6 of 214 counties have no coverage based on EIA data. All six of these counties are in the North weather zone, which has 44 counties; I assign all six 1/44 of North weather zone demand and renormalize the other 38 county weights such that all weights add to one.} The EIA data reports, for the calendar year of 2020, the total delivered electricity by sector (residential, commercial, industrial, and transportation) for each transmission and distribution utility that operated in ERCOT. I allocate those deliveries to counties based on a two-step weighting procedure. First, I look at a county's share of households (per the American Community Survey), and commercial and industrial employment (per the Census Bureau's County Business Patterns) within a utility's service area.\footnote{EIA also reports a small ($<$0.1\%) share of electricity delivered for transportation purposes, which I allocate to counties based on their commercial employment share.} This yields the share of a utility's deliveries that are estimated to belong to each utility-county-sector pair. Then, I generate county weights based on each county's total utility deliveries across utilities. As a result, the EIA data determine the spatial allocation of demand within weather zones, while the ERCOT data determine the level of demand for each hour in the data. 

Table \ref{tab:production-summary} shows summary statistics for the production panel. Average prices increase over the 2018-2025 period; very high prices in 2021 are due to Winter Storm Uri, which hit ERCOT in February and caused widespread chaos throughout the grid, with several important transmission facilities failing, renewables unable to deliver electricity to consumers, and gas lines freezing. Prices were as high as \$10,000/MWh during the storm. Generation and demand increased almost 25\% over the period, with most of the additions coming in the 2022-2024 range. Curtailment increased over the period as well, as more renewables competed for a roughly fixed amount of transmission capacity. Regionally, the Far West had the most growth, with demand more than doubling between 2018 and 2024. Demand is otherwise concentrated in load centers like the Coast, North Central, and South Central. Generation is more evenly spread out, and renewable curtailment is most frequent in regions where generation outstrips demand.\footnote{Curtailment in the Eastern weather zone in 2020 is 12.7\%, which comes from a single solar resource that entered the grid in November of 2020.}

\subsubsection{Infrastructure Panel}
\label{sec:infrastructure}
The infrastructure panel is based on transmission line data, which comes from the Homeland Infrastructure Foundation Level Data (HIFLD). It is a database of transmission lines, ranging in voltage from 69 kV to 765 kV, available in the U.S. It contains information on the geographic location of transmission lines and the voltage of those lines. Transmission line capacity is proportional to voltage. To focus on bulk transmission of electricity, I restrict attention to ERCOT's 230 and 345 kV transmission lines, which are ERCOT's highest capacity lines and most frequently used for bulk transmission.\footnote{In two instances, I found that ERCOT described transmission lines as 345 kV lines while the HIFLD characterized them as 138 kV lines. I assumed that ERCOT's characterization was correct in this case.} I link transmission line data to line-level congestion information available from ERCOT (NP6-86-CD). ERCOT's congestion data reports, at the 5-minute level and between substation pairs, where congestion arose on the grid. It also reveals the transmission capacity in MW between the substations where the congestion happened and the shadow price of that congestion, which I describe in Section XXX. Substations are referred to using standard ERCOT abbreviations.

226 substations appeared in the congestion data for congested elements on bulk transmission lines. Linking ERCOT's congestion data to physical grid segments required a manual review of ERCOT substation names. I used ERCOT's Major Transmission Element (MTE) datasets to map ERCOT's abbreviated station names to full station names. There were several situations where substation names did not appear in the MTE data; in these cases, I performed review of other ERCOT and public reports to map the abbreviations to full substation names. I then matched full substation names to data from OpenStreetMap, a public infrastructure database that contains geolocations for electrical infrastructure worldwide. OpenStreetMap is frequently used in energy research (see, e.g., xxx, yyy, zzz). OpenStreetMap data provided precise locations for correctly mapped substations. For substations that did not have a corresponding entry in OpenStreetMap, I searched documentary evidence to determine the county that they were located in. I was able to successfully find county information for 215 of 226 total substations in the congestion data. 

ERCOT also has a number of Generic Transmission Constraints (GTCs). GTCs monitor a collection of elements and control flows across all of them simultaneously. I identified 26 GTCs in the congestion data, all of which I was able to construct constituent transmission lines or geographic monitoring regions for.\footnote{Region-based GTCs monitor all of the flows into and out of regions. None of the region-based GTCs were used in the model of the grid described in Section XXX because they did not bind on bulk transmission lines sufficiently often.} 

Finally, I used the county-level substation mappings to identify where congestion arose at the county pair level. For the remainder of the analysis, I condense the bulk transmission kV grid in ERCOT into a graph, where counties are nodes and any transmission line that runs through two counties is treated as an edge. I assume that electricity can be stepped down from the bulk transmission grid in any county that has a transmission line running through it. In many situations, multiple transmission lines connect two counties. I recover the implied joint capacity for these transmission lines. Implicitly, I treat intracounty electricity transmission as freely available. A limited number of counties do not have any 230 or 345 kV transmission lines. I assume generation from those counties is freely transmitted to the bulk transmission backbone in proportion to the total number of 138 kV lines that connect the islanded county to counties that lie on the transmission backbone. In total, there are 204 county nodes connected by 300 edges. 10 nodes do not have any 230 or 345 kV transmission lines. 

\subsubsection{Fixed Cost Panel}
\label{sec:fixed_costs}

I construct a panel of yearly fixed cost shifters for both wind and solar generators. I treat capital costs, land costs, and state and federal taxes as fixed cost shifters, which are all either time- or region-varying. I cannot measure some things that clearly contribute to fixed costs, like interconnection fees or labor costs. This omitted variation is not separately identified and loads into the regional fixed-cost intercept that I estimate, described in Section XXX.

I use capital costs from the Lawrence Berkeley National Laboratory 2025 wind and solar market reports (Utility-Scale Solar Update and Land-Based Wind Energy Technology Update). I calculate the price of land based on the Texas A\&M Real Estate Research Center's annual large-tract land price database. The database contains land prices for 33 Land Market Areas, which I mapped to counties by hand. Next, I measure the Texas Chapter 313 policy tax abatement that generators were eligible to enter into with local school districts. The Chapter 313 abatement was a property tax incentive for capital investments, where investor property taxes were capped for school district Maintenance and Operations taxes purposes. Maintenance and Operations taxes vary at the school district level. The program expired at the end of 2022, though abatements negotiated before 2022 were allowed to continue until their original expiry date occurred. I generate an annual value of this program for each project with a negotiated agreement. I apply the Production Tax Credit (PTC) to wind generators and the Investment Tax Credit (ITC) to solar generators. The PTC applies to wind production; I annualize PTC credits to generators based on county-level wind capacity factors and sum PTC credits over a 10-year horizon.\footnote{The next time I estimate the model I will shift this to a net present valuation.} PTC credit rates varied over the 2018-2025 period: in 2018, the rate was \$0.018/kWh; in 2019, \$0.012/kWh; in 2020-2021, \$0.018/kWh; and in 2022-2024, \$0.03/kWh. The ITC is calculated as the ITC rate times the capital cost of solar projects in a given year. The ITC rate over this period was 30\% in all years except 2020 and 2021, when the ITC rate was 26\%. 


\section{Prices and Quantities Under Arbitrary Grid Configurations}
\label{sec:ps_and_qs}

Entry to the electric grid is a function of expected prices and expected quantities. Prices and quantities, in turn, depend on the dispatch process: where electricity comes from and where it travels to. This section details the steps I take to estimate dispatch under any potential configuration of generation on the grid. 

I describe this process in five steps. First, I give a brief overview of the network representation I use and my estimation of power flows.\footnote{Note that my power flow estimation is intended to recover economically relevant congestion exposure and renewable revenues and makes assumptions that would be inappropriate for a full engineering model of the grid.} Next, I describe the dispatch problem in more detail and how it maps to price formation at the nodal level. I then describe the two estimation tasks I have: estimating unobserved grid capacity and modeling the process by which the grid operator commits dispatchable units to produce. Finally, I discuss validation of the model I build. 

\subsection{Network}
\label{sec:network}

The representation of the grid that I use in this article treats counties as nodes connected by edges, as described in Section \ref{sec:infrastructure}. Edges are based on the bulk transmission grid; when any two neighboring counties have at least one 230 kV or 345 kV transmission line connecting them, they are given an edge on the graph representation I use. In cases where multiple transmission lines connect two counties, I aggregate those transmission lines into a single edge. 

ERCOT uses AC transmission. This means that grid operators cannot control where power travels on the network. Instead, grid operators turn generators on and allow demand to take power off the grid at places where its needed. This lack of control means that grid operators must be careful to not turn generators on in such a way that the flow of power will overload any transmission line on the grid. I now describe the process by which grid operators measure power flows on the network.

Power flows depend on net injections and a transmission line's reactance (i.e., how strongly a line limits power flows) to power. Injections matter based on direction and size: nodes that inject electricity to the grid tend to push power towards nodes that withdraw electricity from the grid; larger injections push more power away from a node and larger withdrawals pull more power towards a node. Reactance depends on edge length and edge voltage. Longer edges have higher reactances, while higher voltage edges have lower reactances. 

I assign each underlying transmission line in the grid a reactance rate of $0.32$ ohms/km.\footnote{$0.32$ is standard because of XXXX. Some edges represent multiple transmission lines. I weight underlying edges in proportion to their length and calculate a weighted average reactance across all constituent lines for multiple line edges.} I calculate reactance for edge $c_{ij}$ (which links nodes $i$ and $j$) as a per km measure based on length $L_c$ and voltage $V_c$:
\begin{align}
    X_c = \frac{0.32 \times L_c \times 100}{V_c^2}
\end{align}
Note that reactance rates are common across all edges. Assigning common reactances omits details about equipment like conductor material and transformers from my model. Further, because the grid connections I consider are almost entirely 345 kV lines, almost all heterogeneity in line reactance comes from line lengths. [Will add cites throughout here, this is all standard electrical engineering to my knowledge]

Power flows on edge $c_{ij}$ can now be solved. I use a DC power-flow model to do solve for flows (DC power-flow models are common in industry; ERCOT uses a DC power-flow model to solve dispatch, see XXX). I make several assumptions in order to estimate the DC power-flow model: I ignore resistance and line losses, which are less concerning at shorter distances (relative to, e.g., longer distance transmission considered by \cite{Hausman2024}); and I assume net injections sum to zero (so the grid is balanced) and phase angle differences are small (so there are not blackouts and equipment still works). Allowing $\theta_i$ to represent a node's phase angle (which measures the net injection of node $i$ relative to the rest of the network), power flows $F_c$ on edge $c_{ij}$ are equal to:
\begin{align}
    F_c = \frac{\theta_i - \theta_j}{X_c}.
\end{align}

Solving the grid operator's problem requires understanding how an injection at node $i$ travels throughout the grid. To do so, I construct a nodal susceptance matrix, where edge susceptance is the inverse of edge reactance: $b_{ij} = \frac{1}{X_c}$. Let $B$ be the nodal susceptance matrix, entries of which give the susceptance of all edges on the grid: $B_{ij} = -b_{ij}$. For unconnected nodes, entries are set to 0. Diagonal entries are the sum of susceptances for all edges connected to a node: $B_{ii} = \sum_{j:(i,j)\in \mathcal{E}} b_{ij}$, where $\mathcal{E}$ contains all edges connected to node $i$. Then, nodal balance is satisfied if $p = B\theta$, where $p$ is the vector of county net injections and $\theta$ is the vector of all phase angles (insert cite). 

Phase angles are unidentified in levels; instead, they are defined in relation to a reference county (called a slack county in the literature). I set Harris county as the slack county and modify the $B$ matrix by removing the Harris row and column, yielding the reduced matrix $\widetilde{B}$.\footnote{The reference-node choice is an accounting normalization: for any balanced vector of net injections, predicted physical flows are invariant to the choice of slack node.} The phase angle of each non-Harris county is then given by the following:
\begin{align}
    \widetilde{\theta} = \widetilde{B}^{-1} \widetilde{p}.
\end{align}
$\widetilde{B}^{-1}$ is the inverse of the reduced $B$ matrix, $\widetilde{p}$ is the vector of non-Harris net injections, and $\widetilde{\theta}$ is the corresponding vector of phase angles. 

Finally, we can calculate Power Transfer Distribution Factors $PTDF_{c,r}$, which describe how a one-MW injection in $r$ changes power flows on edge $c_{ij}$ (conditional on a one-MW withdrawal in Harris county). The PTDF for corridor $c_{i,j}$ from an injection in region $r$ is given by the following:
\begin{align}
    PTDF_{c,r} = \frac{1}{X_c}[(\widetilde{B}^{-1})_{i,r} - (\widetilde{B}^{-1})_{j,r}]
\end{align}
Note that injections in region $r$ can have effects on power flows far from the injection site. $c$ does not have to touch $r$. PTDFs are bound by $[-1,1]$, where negative numbers indicate power flowing in the opposite direction from a randomly chosen reference direction for a given corridor.\footnote{A PTDF is bound by $[-1,1]$ because the injection is one MW.} Collecting all PTDFs in a matrix $PTDF$ shows how a one-MW injection in any region changes power flows throughout the grid.

The $PTDF$ matrix can be used to calculate flows on corridors across the grid. Given a vector of net injections $p$, the flows on edge $c$ are given by:
\begin{align}
    F_c = \sum_r PTDF_{c,r} p_r.
\end{align}
Note that this construction of flows, which is what I use in estimation, does not rely on the assumption I made about reactance.

\subsection{Dispatch Problem}
\label{sec:dispatch_problem}
I now discuss the grid operator's problem. Let regions be denoted by $r$, dispatch periods denoted by $t$, and renewable technology denoted by $\tau$. Region-technology-period production is given by $q_{r\tau t}$, total regional production by $q_{rt} \equiv \sum_\tau q_{r\tau t}$, total cost by $\mathcal{C}_{r\tau t}(q_{r\tau t})$, and marginal cost by $MC_{r\tau t}(q_{r\tau t}) \equiv \mathcal{C}'_{r\tau t}(q_{r\tau t})$. Regional demand is given by $D_{rt}$ and system demand aggregates regional demand across regions $D_t = \sum_r D_{rt}$. Let flows on transmission lines be given by $F_{ct}$, with transmission line capacities of $K_c$. 

Grid operators solve a least cost dispatch optimization problem. For a given $t$, they choose generation of each technology (including non-renewable technologies) at each node on the grid to minimize costs. They face two constraints: system demand equals system generation, and flows on transmission lines do not exceed transmission line capacities. The grid operator's problem can be formulated as a Lagrangian:
\begin{equation}
\begin{aligned}
    \mathcal{L}_t
    ={}& \sum_r \sum_\tau \mathcal{C}_{r\tau t}(q_{r\tau t})
    + \lambda_t \left(\sum_r D_{rt}
        - \sum_r \sum_\tau q_{r\tau t}\right) \\
    &+ \sum_c \mu^{c,+}_t \left(\sum_r PTDF_{c,r}
        (q_{rt} - D_{rt}) - K_c\right) \\
    &+ \sum_c \mu^{c,-}_t \left(-\sum_r PTDF_{c,r}
        (q_{rt} - D_{rt}) - K_c\right).
\end{aligned}
\end{equation}
$\lambda_t$ is a Lagrangian multiplier on the system demand and $\mu^{c,+}_t$ and $\mu^{c,-}_t$ are Lagrangian multipliers on transmission constraints. $\mu^{c,+}_t$ and $\mu^{c,-}_t$ are often referred to as the shadow price of transmission capacity, indicating the value of relaxing transmission constraints by one MW.\footnote{Congestion costs, as discussed in Section \ref{sec:policy}, are calculated as $\sum_t\sum_c[\mu^{c,+}_t F_{ct}+\mu^{c,-}_t(-F_{ct})]\delta t$, summing the shadow price of the constraint over the period of time ($\delta t$) and in the direction in which it was constrained.} Note that line constraints have both upper and lower limits, where sign depends on the direction in which power is flowing on the transmission line. For clarity, I omit generator production capacity limits.

Let the grid be uncongested, so $\mu^{c,+}_t=0$ and $\mu^{c,-}_t=0$ across all grid edges. Consider a marginal generator whose production is not at full capacity. Taking the derivative with respect to $q_{r\tau t}$ yields the system price of electricity:
\begin{align}
    \frac{\partial \mathcal{L}_t}{\partial q_{r\tau t}} = MC_{r\tau t}(q_{r\tau t}) - \lambda_t = 0.
\end{align}
Let $\mathcal{C}^\star_t$ be the minimized system dispatch cost and $D_t \equiv \sum_r D_{rt}$. Then the price of electricity in an uncongested system is set by the marginal cost of the marginal generator. By the envelope theorem, the cost of increasing demand by one MW is $\lambda_t$:
\begin{align}
    \lambda_t = \frac{\partial \mathcal{C}_t^\star}{\partial D_t}.
\end{align}

As the grid becomes congested, $\mu$s turn on and local prices diverge from system prices. Taking the derivative with respect to $q_{r\tau t}$ yields the price of electricity in region $r$:
\begin{align}
    \frac{\partial \mathcal{L}_t}{\partial q_{r\tau t}}
    = MC_{r\tau t}(q_{r\tau t}) - \lambda_t
    + \sum_c (\mu^{c,+}_t-\mu^{c,-}_t)PTDF_{c,r} = 0.
\end{align}
Rearranging gives:
\begin{align}
    LMP_{rt}
    = \frac{\partial \mathcal{C}_t^\star}{\partial D_{rt}}
    = \lambda_t-\sum_c(\mu^{c,+}_t-\mu^{c,-}_t)PTDF_{c,r}
    = MC_{r\tau t}(q_{r\tau t}).
\end{align}
$LMP_{rt}$ is called the locational marginal price for a given region. 

The dispatch problem consists of mostly data: $\mathcal{C}_{r\tau t}(q_{r\tau t}), q_{r\tau t}, \text{ and } \lambda_t$ are directly observable in ERCOT datasets; $D_{rt}$ is estimated as described in Section \ref{sec:production}; and $PTDF_{c,r}$s are parameters. When congestion arises on the grid, ERCOT reports the shadow price of that congestion, giving evidence for what $\mu^{c,+}_t$ and $\mu^{c,-}_t$ were in relevant periods.\footnote{Note, however, that $q_{r\tau t}, \lambda_t, \text{ and } \mu^{c,+}_t \text{ and } \mu^{c,-}_t$ will later be endogenous outputs of the counterfactual dispatch model.}

Data on $K_c$ is unavailable. I discuss how I estimate $K_c$ in the following section. 

\subsection{Grid Capacity}
\label{sec:grid_capacity}

After reducing the network to the county-edge graph, I have 300 total edges for which I need capacity information. I estimate $K_c$ in four parts. First, I set initial capacities for each edge in the graph. Initial capacities are based on direct limits taken from the shadow price data when available. The shadow price data only provides limits for transmission lines that are actively constrained. Of the 300 nodes, 50 receive estimated limits based on the shadow price data. The remaining 250 edges receive capacities based on their voltage and the number of lines aggregated within a single edge. 

Voltage measures the pressure that transmission exerts on power to flow in a given direction -- not the maximum amount of electricity that flows along a given transmission line. The two are proportional however. I assume that all remaining 345 kV edges have a capacity of 1200 MW, which is the average binding limit for 345 kV transmission lines based on the congestion price data. For consistency, I assume that 230 kV transmission lines have a capacity of 800 MW.\footnote{A 1,200 MW real-power limit on a balanced three-phase 345 kV AC circuit corresponds to a line current of about 2.0 kA:
\[
I = \frac{1200\text{ MW}}{\sqrt{3}\times 345\text{ kV}} \approx 2.01\text{ kA},
\]
where $I$ is line current. I interpret this value as a short-corridor thermal-capacity prior rather than a long-distance transfer capability. This distinction is consistent with engineering loadability guidance: AEP reports that 345 kV lines have much lower long-distance loadability under the surge-impedance-loading yardstick, but that short-distance loadability can be higher when conductor and station-equipment thermal ratings are adequate; it gives 1,500 MW on a 345 kV line as approximately 3.8 SIL over about 50 miles (\cite{AEPTransmissionFacts,DunlopGutmanMarchenko1979}). Consistent with NERC facility-rating standards, observed facility ratings depend on equipment ratings, ambient assumptions, operating limits, and verified engineering analysis (\cite{NERCFAC0085}). There are 25 edges with voltage classes of 230 kV. Assuming these edges have a capacity of 800 MW is consistent for a 230 kV line with a line current of about 2.0 kA. In reality, capacity depends on many factors, including conductor size, the number of conductors, weather, substation equipment, line length, and more.} 

In the second step of the process, I use evidence from observed dispatch. Based on observed net injections in a county -- i.e., looking at, for each period from 2018-2024, observed generation and estimated demand -- PTDFs allow me to calculate the implied flows on a given edge as follows:
\begin{align}
    F_{ct} = \sum_r PTDF_{cr}(q_{rt} - D_{rt}).
\end{align}
If an edge has direct evidence of an implied power flow greater than its initial capacity and it did not cause congestion during period $t$, I adjust the capacity of that line upwards to fit that data. 

In the third step of the process, I solve the grid operator's dispatch problem for each hour from 2018-2024. To do so, I reconstruct historical supply stacks using generator offers from the SCED data and observed dispatch limits.\footnote{I also account for committed units required output.} Then I solve dispatch based on the current best guess of edge capacities. 

I compare the solved dispatch model to observed outputs at the edge level. For each edge, I look at the spread between the model implied and observed nodal endpoint prices. I calculate 90$^{th}$ percentile spreads; if the model spread is too high, I increase edge capacity, and if it's too low, then I decrease edge capacity.\footnote{I look at 90$^{th}$ percentile spreads because differences are only informative when constraints are binding. Maximum spreads measure prices during blackouts, which are difficult to accurately price.} Then I look at the modeled vs observed binding frequency for each edge. I define observed binding as anytime the edge's absolute endpoint price spread is greater than \$5 MWh. \footnote{This part of the process is not that important but the way I'm doing this is pretty ugly, so worth reconsidering for the next time I run the model.} 

This process results in small adjustments to modeled capacities. It increases flows for 3\% of edges and decreases flows for 17\% of edges. The maximum decrease is 356 MW, the 10$^{th}$ percentile decrease is 26 MW, and the 25$^{th}$ percentile decrease is 0 MW. Similarly, the maximum increase is 391 MW, while the 99$^{th}$ percentile increase is 61 MW and the 95$^{th}$ percentile increase is 0 MW. It improves the prices calculated by the dispatch model relative to the observed data by roughly 10\%. 

[to come: brief description of how I handle the WESTEX GTC. deep appendix material.]

\subsection{Commitment}
\label{sec:commitment}

This section describes the steps I take to account for unit commitment in the dispatch process. As shown in Section \ref{sec:grid}, unit commitment places a (sometimes large) stack of price insensitive dispatchable supply at the bottom of the offer stack. This shifts the offer curve for generators bidding into the market based on marginal cost to the right and depresses prices as a result. In order to have an accurate counterfactual model of prices and quantities, I must account for the grid operator's commitment behavior in equilibrium. 

A model of endogenous commitment is outside the scope of this paper. Instead, I estimate a reduced form approximation of commitment each period. Estimating commitment requires two steps. First, I estimate the MW capacity of dispatchable resources available to commit each hour. I estimate availability separately for coal and natural gas. The model depends on unit online statuses. Online statuses reflect whether or not a generator is currently producing electricity. Online dispatchable resources must produce some electricity (or else they pay ramping costs to turn off). They are able to produce up to a High Sustained Limit (HSL). I model the expected amount of available coal as equal to the sum of HSLs for all online coal units in a given hour.\footnote{This assumes that counterfactual entry would not change whether or not units were online in any given period.} I model the expected amount of available natural gas as the sum of HSL across online units, plus $\alpha$ times the HSL of all offline gas units. I find that setting $\alpha=0.40$ produces the best fit with prices. Implicitly, this suggests that as dispatch conditions shift, 40\% of offline gas is available for commitment. 

The second step of the commitment model estimates actual commitment. To estimate the model, I predict commitment separately by county and fuel based on local and system net load, renewable production, and system tightness, which measures how close system net load is to the total HSL of all online and offline coal and natural gas units. I include fixed effects for hour, month, and market-period.\footnote{I include five market periods in the model: pre-Winter Storm Uri, from January 1, 2018 to February 12, 2021; Uri, from February 13 to 19, 2021; Post-Uri, from February 20 to June 7, 2021; 2021 Reform, from June 8 to December 31, 2021; and Post-reform, from January 1, 2022 onwards. ERCOT adopted a more conservative operating posture starting in summer of 2021, committing resources more actively. In 2022, ERCOT made changes to its Operating Research Demand Curve, which determines how the units that ERCOT forcibly commits get compensated for that commitment.} The regression produces, for each counterfactual hour, dispatchable availability, which I apply to submitted coal and natural gas offer curves. 

\subsection{Validation}
\label{sec:pq_validation}
To test the fit of my capacity and commitment models, I estimated the commitment rule using data from 2018-2022. I then evaluated the prices that the model gave for 2023 against observed price data in 2023:
\begin{table}[!htbp]
\centering
\caption{Out-of-Sample Validation of the Commitment and Dispatch Model}
\label{tab:commitment_validation}
\begin{threeparttable}

\begin{tabular}{lrrrrrr}
\toprule
& Model & Observed & Bias & MAE & RMSE & Correlation \\
\midrule
\multicolumn{7}{l}{\textit{Panel A: System price validation}} \\
Demand-weighted LMP (\$/MWh)
    & 52.38 & 52.24 & 0.14 & 29.14 & 192.30 & 0.85 \\
\addlinespace
Load-shedding frequency
    & 0.00 & --- & --- & --- & --- & --- \\
\midrule
\multicolumn{7}{l}{\textit{Panel B: Renewable price-value validation}} \\
Solar capture price (\$/MWh)
    & 71.89 & 67.67 & 4.22 & --- & --- & --- \\
Wind capture price (\$/MWh)
    & 25.26 & 22.98 & 2.28 & --- & --- & --- \\
Solar revenue (\$/MW-year)
    & 116{,}078 & 109{,}775 & 6{,}303 & 13{,}715 & 22{,}089 & 0.86 \\
Wind revenue (\$/MW-year)
    & 64{,}378 & 58{,}567 & 5{,}811 & 13{,}674 & 19{,}285 & 0.61 \\
\bottomrule
\end{tabular}

\begin{tablenotes}[flushleft]
\footnotesize
\item \textit{Notes:}
The table evaluates the commitment and network-dispatch model on 2,000
evenly sampled hours from 2023, which are excluded from estimation of the
commitment-rule coefficients. Model prices are generated using the fixed
$\alpha=0.40$ expected-online-stack rule and the contemporaneous post-reform
commitment regime. Panel A uses the 1,999 sampled hours with matched observed
county LMPs. The model sheds no load in the validation sample. In Panel B,
observed renewable output is held fixed and valued at modeled and observed
county LMPs. Revenue statistics are annualized by weighting each sampled hour
by $8{,}760/2{,}000$. The solar and wind statistics average county-level
per-MW outcomes across 23 and 44 electrical county nodes, respectively.
RMSE is sensitive to a small number of scarcity-price observations.
\end{tablenotes}

\end{threeparttable}
\end{table}

The model shows very strong fit at the system price and renewable price levels. Model demand-weighted LMPs in 2023 are \$52.38, while observed demand-weighted LMPs are \$52.24. The Mean Absolute Error between model and observed prices is higher -- \$29.14 -- though much of that discrepancy is due to tail events driving up mean. Median absolute error, for example, is \$3.89. The correlation across prices is also very high -- 85\%. Note that RMSE is high because of tail events, which are difficult to price accurately. Tail events are important for accurately estimating entrant revenues, as they have the potential to increase revenue dramatically. The difficulties pinning down tail revenues do not cause issues for expected wind and solar yearly per-MW revenues, which are quite accurate. 

\section{Stationary Entry Model}
\label{sec:entry_model}
This section details the empirical model of generator entry. I use to examine entry patterns on the Texas grid. I treat generators as technology-specific, independently operated, single-plant firms.\footnote{Insert description of joint ownership in Texas, \cite{JohnstonLiuYang2024} find XXX.} I fix entrant sizes to the average nameplate MW of wind and solar entrants from 2018-2024. Each period, generators make a single decision in the model: enter or permanently forgo the opportunity. Next, generators exit. Generator exit is rare, which I capture with an exogenous and reduced form process.\footnote{Five wind resources exited during the 2018-2024 period, while zero solar resources did.} In the baseline version of the model, I omit transmission line upgrades, reflecting the limited number of bulk transmission grid improvements ERCOT made over this period. I return to the upgrade process in counterfactuals. The final stage of the model is production, where the grid operator dispatches supply to meet demand, using the dispatch model described in Section \ref{sec:dispatch_problem}. 

I abstract away from interconnection queue considerations by modeling entry decisions as they occur at the end of the interconnection queue (see, e.g., \cite{JohnstonLiuYang2024}, \cite{Boatwright2025} for a discussion of the queue). This reflects the more straightforward approach to the queue in ERCOT: there are no generator-funded network upgrades, so the timing of entry does not determine the fees that generators face upon entry; and the queue is shorter (projects wait in the queue for an average of two years) than it is elsewhere (\cite{LBNL2024}). Because the decision is modeled at the end of the queue, information revealed during the interconnection study is reflected in the entrant's state at the time of the decision.

Entrants face site-specific entry opportunities, so the entry decision is at the region-technology level. Regions in the model are county clusters -- there are 40 regions in total and each region has on average six counties within it.\footnote{Regions are generated based on edge characteristics, assessed on distance (smaller distances are better), capacity similarity (high-capacity corridors are grouped together), and congestion frequency (congestion separates regions). The goal is to cluster counties with limited intraregional congestion into a larger region. Implicitly, this says that when generators site in that region, intraregional deliverability is more robust than cross-region deliverability.} The number of generators eligible to enter each region each period, $\bar{N}_{r\tau}$, is fixed at the largest number of region-year entrants seen in each region across all years plus one. Solar and wind generators made up the vast majority of entrants to the grid over the period I study, though battery storage entered in large quantities in 2023 and 2024. I treat all incumbents and new non-solar or wind entrants as exogenous shifts in supply. 

\subsection{Timing and Information}

Each period represents one year. The timing of the model proceeds as follows:
\begin{enumerate}
    \item States and beliefs are realized. Entrants draw T1EV i.i.d. cost shocks and simultaneously decide on entry. 
    \item Retirement is realized. If chosen for retirement, generators exit. Generator entry occurs.
    \item Production occurs for all studied hours using transmission capacity and supply curves constructed based on regional technology quantities. Payoffs accrue to generators and states transition.
\end{enumerate}

Publicly observed state variables are given by $x_{r\tau}$:
\begin{align*}
   x_{r\tau} = \{r, \tau, q_{r,\mathrm{solar}}, q_{r,\mathrm{wind}}, Q_{\mathrm{solar}}, Q_{\mathrm{wind}} \},
\end{align*}
where $r$ indicates the region in which an entrant is located, $\tau$ indicates the entrant's technology type, $q_{r,\mathrm{solar}}$ indicates the quantity in MW of solar generation in $r$, $q_{r,\mathrm{wind}}$ indicates the quantity in MW of wind generation in $r$, and $Q_\tau$ indicates the total quantity in MW of each technology type on the grid. I omit the subscripts on $x$ moving forward for clarity.

Forecasting payoffs requires more than the observed states. Consider, for example, $q_{r,\mathrm{solar}} << Q_{\mathrm{solar}}$. If solar generation outside region $r$ was uniform across all other regions, congestion would be low in region $r$. If outside solar generation was instead concentrated in a single region, congestion would be rampant on the grid, and prices in $r$ would either be very low (if it was further from demand than the concentrated region) or very high (if it was closer to demand than the concentrated region). These situations, too, are not equally likely: other-region entry concentrating in a single region would be uneconomic. Entrants have beliefs over $g$, which is the full county-level allocation of generation. The belief over entry allocations $M(g | x)$ provides the distribution of $g$ that is consistent with $x$ for each $x$ (\cite{IfrachWeintraub2017}).\footnote{In equilibrium, $M$ must coincide with the equilibrium entry allocation across regions.} In effect, entrants have a distribution of potential other region quantities for each $Q$ that corresponds to how often that allocation of quantities would occur under the solved entry policy $p(x)$ and retirement. 

Flow payoffs depend on $g$, which determines the LMPs in the model. Entrants do not observe the particular allocation $g$ when choosing whether to enter. Instead, they evaluate payoffs and transitions under the equilibrium conditional belief $M(g\mid x)$. The remaining beliefs depend on $p(x)$: $q$ transitions according to realized entry counts and retirement. $Q$ evolves under entry implied by $p(x)$.

I now describe the model, starting with production.

\subsection{Production}

Flow payoffs are based on expected revenues from the dispatch model. The dispatch model, as described in Section \ref{sec:ps_and_qs}, solves for least-cost production sufficient to meet demand. The system cost of electricity is given by $\lambda_t$, which reflects what the price of electricity would be absent any transmission congestion. A region's price is given by 
\begin{align}
    LMP_{rt} = \lambda_t + \text{congestion costs}_{rt}.
\end{align} Regions on the demand-side of congestion have positive congestion costs, while regions on the generation-side of congestion have negative congestion costs. 

Let $\bar q_\tau^{\mathrm{ent}}$ denote the capacity of a technology-$\tau$ entrant, let $\omega_t$ denote the annual weight assigned to dispatch hour $t$, and let $q^p_{r\tau t}(x,g)$ denote the entrant's dispatched production. Its flow payoff is
\begin{align}
    \pi_{r\tau}(x,g)
    = \sum_t \omega_t q^p_{r\tau t}(x,g)LMP_{rt}(x,g),
\end{align}
Note that resource quality is embedded within $q^p_{r\tau t}$. If we let $q_{r\tau}$ indicate region-technology capacity and $f_{r\tau} \in (0,1)$ indicate a region-technology pair's capacity factor, $q^p_{r\tau} \leq f_{r\tau}q_{r\tau}$. I use $\pi_\tau$ as a shorthand for $\pi_{r\tau}(x,g)$ moving forwards. 

Expected payoffs integrate over $g$
\begin{align}
    E_g[\pi] = \sum_g M(g|x) \pi(x,g).
\end{align} 

Hours are defined by the capacity factor of renewables (solar does not produce at night, for example) and the level of demand that must be met. Instead of simulating payoffs for all 8,760 hours in a year, I simulate payoffs for a subset of hours. The chosen subset of hours is representative across capacity factors, demand, demand geography, time of year, and time of day. Current iterations of the model are based on 96 sample hours. Payoffs are then constructed as a weighted sum of payoffs in each representative hour, with weights corresponding to the frequency of an hour's conditions occurring over the course of the year. 

The model is infinite-horizon and stationary. Demand in ERCOT rose by roughly 2\% per year between 2018 and 2024. Sampling hours across years thus risks asking 2018-level capacity to meet 2024-level demand, which would lead to blackouts. Instead of doing this, I sample hours exclusively from 2018. To the extent that entry happened in response to the observed increase in demand, my model will attribute that entry to decreasing fixed costs. There were also regional shifts in demand in ERCOT over this period. Oilfield electrification in the Far West led to localized demand doubling in the region and increased congestion. I test the effect of this demand increase by doing XXX [and show it isn't important?]

\subsection{Exit}
\label{sec:exit}
% TODO: give these (and other objects) greek symbols and report the values I use in a table for the results section.
I assign each generator an exit rate based on their technology type. Solar generators retire with probability 1/30 each year and wind generators retire with probability 1/25 each year. These hazards imply operating lives of 30 and 25 years, respectively, which correspond to the average lifespan of utility-scale wind and solar technology. I do not treat retirement as increasingly likely with generator age, primarily for computational reasons.\footnote{An age-independent geometric lifetime is appealing because I do not have to keep track of generator ages.} 

This is primarily a technical assumption. Retirement is rare in sample -- only five wind generators retired over the period I study -- and arrives randomly. The shortest lived wind generator in ERCOT exited after 11 years of production, for example. Wind and solar marginal costs are zero in the model and wind generators often report negative bids (see \ref{sec:grid}); modeling strategic retirement would be challenging. The retirement rates I use are sufficiently low that retirement is unlikely for any individual generator over the time span I look at. Additionally, adding retirement into the model prevents generation from accumulating indefinitely and makes the stationary distribution of generation less reliant on the model's initial conditions. 

\subsection{Entry}

Each period $t$, a new set of entry opportunities $\bar{N}_{r\tau}$ appears in each region. Entrants decide whether they would like to pursue the entry opportunity $a_{r\tau}=1$ or let it lapse $a_{r\tau}=0$. Entry opportunities within a region-technology type are symmetric, thus all entrants play the same mixed strategy entry policy $p(x)$. Entrants must pay a fixed cost $F_{r\tau}$ to enter. 

Entrants are forward-looking and make decisions about entry based on flow payoffs along with future grid evolution, implicitly trading off the resource quality of a location against its current and future expected transmission congestion. They observe the public moment state $x_{r\tau} = \{r, \tau, q_{r,\mathrm{solar}}, q_{r,\mathrm{wind}}, Q_{\mathrm{solar}}, Q_{\mathrm{wind}}\}$ and a vector of private shocks $\{\epsilon_{0r\tau}, \epsilon_{1r\tau}\}$, which I assume follow the T1EV distribution with scale parameter $\sigma_\tau$. They do not condition their entry decision on the full other-region allocation $g$; instead, they hold equilibrium beliefs $M(g\mid x)$ over the allocations consistent with the observed moment state.

Let $\beta$ be the annual discount factor, common to entrants and incumbents. Let $x^\prime$ indicate next period's state realization, $E^{inc}_{p,M}$ indicate the expectation for an incumbent over entry policy $p$ and capacity allocation belief $M$ and $E^{ent}_{p,M}$ indicate the expectation for an entrant. Let $\pi_\tau$ denote the technology-specific flow payoff for a given $x$ and $g$ realization. Then, for a given entry policy $p$ and allocation belief $M$, the value of a generator of type $\tau$ on the grid $V_\tau(x)$ and the value of entry $v_{\tau,1}(x)$\footnote{Note that, for computational reasons, I include a ceiling on entry in each region, which I describe below. This formulation of entry values omits the risk of an entrant being rejected due to a region reaching its capacity ceiling.} are given by:
\begin{align}
    \label{eq:value_func}
    V_\tau(x) &= E^{inc}_{p,M} [\pi_\tau + \beta V_\tau(x^\prime)|x] \\
    v_{\tau,1}(x) &= E^{ent}_{p,M}[\pi_\tau + \beta V_\tau(x^\prime)| x, a_{r\tau}=1]
\end{align}
For a generator of technology $\tau$ already on the grid, the value of being on the grid depends on the generator's production profile, which depends on $\tau$, rival generator entry in all regions, given by $p$, and the generator's belief about how capacity is allocated across other regions, given by $M(g|x)$. Then, retirement is realized and dispatch is run, generating $\pi_\tau$. Finally, states transition. The value of entering is almost identical, except the entrant needs to account for how its own entry affects state variables $q_{r\tau} \text{ and } Q_\tau$ and dispatch realizations.

Entrants weight the value of entry against fixed costs and cost shock realizations:
\begin{align}
    \tilde{v}_{\tau,1}(x) &= v_{\tau, 1}(x) - F_{r\tau} - \epsilon_{1r\tau} \\
    \tilde{v}_{\tau,0}(x) &= - \epsilon_{0r\tau}
\end{align}
Let $p_\tau(x)$ denote the entry probability for a technology-$\tau$ opportunity in state $x$, and let $p \equiv \{p_\tau\}_{\tau\in\mathcal{T}}$ denote the complete entry-policy profile. Projects enter when $\tilde{v}_{\tau,1}(x) \geq \tilde{v}_{\tau,0}(x)$. T1EV shocks yield the logit entry probability:
\begin{align}
    \label{eq:entry_policy}
    p_\tau(x)
    = \frac{\exp\left(\left[v_{\tau,1}(x)-F_{r\tau}\right]/\sigma_\tau\right)}
    {1+\exp\left(\left[v_{\tau,1}(x)-F_{r\tau}\right]/\sigma_\tau\right)}.
\end{align}
$\sigma_\tau$ controls the scale of the cost shock that entrants receive. Small $\sigma_\tau$ values indicate that entry is fairly deterministic conditional on the modeled value of entering in a region, while large $\sigma_\tau$ values indicate substantial unobserved heterogeneity in entry costs. I describe the process I use to estimate $\sigma_\tau$ for each technology in Section \ref{sec:estimation}.

% TODO: describe the below in estimation and identification
%I test a grid of $\sigma$s ranging from 4 to 4,096 for each technology. I use a test and train metholodogy to identify $\sigma$ values that best fit the data. I 
%I fit entry using $\sigma$ to all regions except one using the modeled entry values. Then I predict entry in the left-out region 
%I find that a $\sigma_{\mathrm{solar}}$ of \$24 million USD / project and a $\sigma_{\mathrm{wind}}$ of \$96 MUSD / project produce entry rates consistent with observed entry levels and locations. $\sigma$ values were assessed using a grid of values ranging from $\sigma=4$ to $\sigma=4,096$.

$p(x)$ gives the entry policy that generators play. Inverting Equation \eqref{eq:entry_policy} yields a closed-form representation of fixed costs:
\begin{align}
    \label{eq:fixed_cost_inversion}
    F_{r\tau} = v_{\tau,1}(x) - \sigma_\tau \log(\frac{\hat{p}_\tau(x)}{1-\hat{p}_\tau(x)}).
\end{align}
$\hat{p}(x)$ is observed entry frequencies. I then parameterize the model-implied fixed costs according to observable fixed cost shifters, including capital costs, tax rebates, land costs, and the Chapter 313 tax rebate discussed in Section \ref{sec:fixed_costs}. I aggregate capital costs, land costs, and tax rebates into a $C^{\text{net capex}}_{r\tau}$ value, while I include other non-price shifters in $Z^\prime_{r\tau}$, which I discuss in more detail in Section \ref{sec:estimation}. Letting $\theta_\tau$ hold the coefficients for all non-price shifters, the parameterization is as follows:
\begin{align}
    \label{eq:fixed_cost_representation}
    F_{r\tau}(\theta_\tau) = C^{\text{net capex}}_{r\tau} + Z^\prime_{r\tau}\theta_\tau.
\end{align}
The coefficients on this cost function are estimated using observed entry choices and model-generated entry values. The coefficient on net capex is fixed to one because it is measured in dollars. Projecting fixed costs onto observables provides a common fixed cost function for counterfactuals.

% TODO: describe the below in estimation and identification
%I iterate over the inversion until model entry policies and fixed costs rationalize each cell's entry. I project these fixed costs against technology-specific observables, including capital expenditure costs, tax rebates, land costs, and a Texas tax break for renewables that varies at the county level and which expired in 2023. 

For computational reasons, I specify $\bar{A}_r$, a regional max renewable capacity absorption ceiling. This ensures that the model cannot add generators to a single region unboundedly. The ceiling is a fixed parameter and does not vary over time. I calculate the ceiling as: 
\begin{align}
    \bar{A}_r = [\max \{D^{peak,2018}_r + E^{2018}_r, \max_{t \in \{2018,\dots,2024\}} (q_{rt,\mathrm{solar}} + q_{rt,\mathrm{wind}})\}].
\end{align}
$D^{peak,2018}_r$ is the highest demand observed in region $r$ during 2018 and $E^{2018}_r$ is the maximum export capacity of region $r$ in 2018. Together, this quantity represents the total capacity for consumption and exporting in region $r$. Because grid capacity is fixed in estimation, this quantity holds across all periods. The second term, $\max_{t \in \{2018,...,2024\}} (q_{rt,\mathrm{solar}} + q_{rt,\mathrm{wind}})$, represents the maximum installed capacity in each region across the sample period. Including this term in $\bar{A}_r$ ensures that observed electricity capacity never exceeds the production ceiling that I implement. Note: in a simplified version of the model, the second term never exceeded the first; I'll ensure that's true in the full model and highlight here if so. The ceiling is intended to be unlikely to bind. 

Wind and solar entrants are jointly counted against the ceiling. On rare occasions, entrants may be barred from joining a region because incumbent plus entrant MW is too large to fit under the ceiling. If this arises during a period where multiple entrants attempt to join a region, I assume that entry is allocated randomly across all entrants for that period such that $q^{post}_{r,\mathrm{solar}} + q^{post}_{r,\mathrm{wind}} \leq \bar{A}_r$, where post indicates post-entry quantity.

I rule out strategic delay in entry. There are several reasons this is appropriate in my setting. First, strategic delay requires the foregone revenue from not entering is smaller than any prospective declines in fixed costs. While fixed costs decreased over the period I study, revenues were also high. [insert yearly revenue, estimated fixed costs]... 

\subsection{Equilibrium}

Equilibrium in this model is driven by the tradeoff between resource quality and dynamic congestion concerns. As congestion becomes more likely and more costly, this tradeoff becomes more acute. 

Capturing these dynamics requires a representation of renewable capacity on the grid and where that capacity is located. With 40 regions, estimating a Markov Perfect Equilibrium for this setting would quickly become infeasible. Following \cite{WeintraubBenkardVanRoy2008} and \cite{IfrachWeintraub2017}, I solve for a stationary moment-based Markov equilibrium (MME) where firms track the spatial allocation of rival capacity. Generators include details about their own region (renewable capacity) and a statewide summary of total renewable capacity according to their technology on the grid ($Q_\tau$). 

A stationary equilibrium in this model yields a triple $(p^\star, V^\star, M^\star)$, where $p^\star$ gives entry probabilities in each region and for each technology and $x$; $V^\star$ gives the value to a generator for being on the grid in each region and for each technology and $x$, and $M^\star$ gives beliefs about the allocation of generation across other regions for each $x$. I require that:
\begin{enumerate}
    \item Given $M^\star$ and rival policy $p^\star$, $V^\star$ satisfies the incumbent Bellman equation, 
    \item Entrant values constructed from $V^\star$ generate $p^\star$ based on the logit best response,
    \item Entry, retirement, and stock transitions under $p^\star$ generate a stationary allocation distribution whose conditional law is $M^\star (g \mid x)$.
\end{enumerate}

Equilibrium computation requires the following steps:
\begin{enumerate}
    \item Initialize $M$ and $p$;
    \item Holding $M$ fixed, solve the $V-p$ fixed point;
    \item Simulate entry and retirement based on $p$;
    \item Update $M(g \mid x)$ based on the simulated allocation of generation across regions;
    \item Update $V$ and $p$ and repeat until policy, values, and the allocation law are stationary.
\end{enumerate}

\section{Estimation}
\label{sec:estimation}

This section describes estimation of the dynamic entry model. The model seeks to construct entry values such that they rationalize observed entry choices at the region-technology-year level. The estimation process is summarized in Figure \ref{fig:estimation_identification}. The relevant objects in this section are as follows: $p$, entry probabilities; $V$, an entrant's (or incumbent's) value of being on the grid; $M(g|x)$, an entrants beliefs over the allocation of renewable capacity across other regions; $F$, the fixed costs of entry; and $\sigma$, the scale of entrant cost shocks. I solve two nested fixed points in order to estimate the model. In the first fixed point, I solve for equilibrium based on flow profits from the dispatch model and entry probabilities from finding a fixed point in $V$-$p$. In the second fixed point, I add in entrant beliefs over $M(g|x)$, which must be accurate in equilibrium. Then, I project fixed cost thresholds onto the fixed cost panel to obtain a representation of fixed costs that I can take to counterfactuals. Doing so requires me to re-estimate the nested fixed point routine. In the fourth step, I evaluate how the model fits for a grid of $\sigma$ values. I suppress region and technology subscripts throughout this section for clarity.

% TODO: where to describe that i solve the entry model for 2018? 

\begin{figure}[htbp]
    \centering
    \includegraphics[width=\textwidth]{../figures/estimation/estimation_identification_nested_fixed_points.pdf}
    \caption{Estimation Flow}
    \label{fig:estimation_identification}
\end{figure}

\subsection{Values and Entry Policy}

Solving for an entry equilibrium in profits requires a set of policies $p$ and other-region entry beliefs $M$ such that the resulting project values $V$ rationalize entry for a given $F$. To solve the profit fixed point, I hold $M$, $F$, and $\sigma$ fixed. The profit fixed point finds a fixed point in equations \ref{eq:value_func} and \ref{eq:entry_policy}. 

I begin by generating $V_0$. To do so, I initialize $p$ and $M$. $p_0$ is initialized by regressing $x$, grid capacity, resource quality, and incumbent revenues on observed entry in each region. The regression predicts the total number of entrants in for each period-region-technology. Results for 2018 are used to generate $p_0$. Initial beliefs $M_0$ are drawn from empirical spatial allocations that arose between 2018 and 2024. Entrants observe $x$, but not the full allocation $g$. Instead, they know how $g$ maps into flow profits and future state transitions, and integrate over possible values of $g$ when evaluating entry. I approximate this expectation by giving each prospective entrant three representative draws from $M_0(g\mid x)$. $p_0$ and $M_0$ determine $\pi_\tau$ via the dispatch model for each region. Project entry values $V$ are the discounted flow of future profits. 

With $V$ in hand, I use \ref{eq:fixed_cost_inversion} to recover $F_0$, the fixed costs implied by the model. I initialize $\sigma_0$ to 8 MUSD/project. $\hat{p}$ indicates empirical entry rates, thus yielding the implied threshold that rationalizes observed entry. With $V$ and $F$ recovered, I use \ref{eq:entry_policy} to update $p_0$ to $p_1$. I then update value functions and iterate until reaching a fixed point in $V^\star$ and $p^\star$.

\subsection{Stationary MME}

The values and entry policy fixed point held $M_0$ fixed. In the equilibrium, $M^\star$ must coincide with the stationary allocation distribution of generation. Because $p^\star$ shifts entry probabilities, it shifts $M^\star$: when equilibrium entry shifts in a given region, beliefs about other-region entry must also shift. 

Estimating the stationary MME applies $p^\star$ to $M_0$, similar to the forward simulation method from \cite{IfrachWeintraub2017}. As discussed, I take three draws from $M_0$ to generate allocations of capacity on the grid. $\bar{N}$ entrants then appear in each region. Each potential entrant receives a cost shock $u$ and enters if $u < p^\star(x)$. 

Following entry decisions (and retirement), $g$ transitions $g^\prime = T(g_{M_0}, u; p_1)$; $q$ transitions $q^\prime = \sum_{c \in r} g^\prime_c$; and $Q$ transitions $Q^\prime = \sum_r q^\prime_r$. This yields $x^\prime$. I update $M_0$ to $M_1$, where $M_1$ contains a dampened combination of original allocations and destination allocations observed following state transitions.\footnote{Note that $M(g|x)$ is conditional on $x$, so only allocations corresponding to the same state variables are included in the same distribution.} Having a new $M$ distribution shifts implied $p$ and $V$ values, thus restarting the values and entry policy loop. I continue estimating $p$, $V$, and $M$ until a stationary MME is reached. 

\subsection{Fixed Costs}
%\textbf{Note: I'm reviewing ways to more convincingly identify $\sigma$. It's been easy for solar, but difficult for wind.}

The stationary MME produces the region-technology-year threshold for fixed costs that aligns model and observed entry via equation \ref{eq:fixed_cost_inversion}. I parameterize the fixed cost threshold according to observable cost shifters using equation \ref{eq:fixed_cost_representation} (reprinted for clarity):
\begin{align*}
    F_{r\tau}(\theta_\tau) = C^{\mathrm{net capex}}_{r\tau} + Z^\prime_{r\tau}\theta_\tau.
\end{align*} 

Using ordinary least squares to parameterize the fixed cost function is not well-behaved because many region-technology-year cells have no entry, meaning the fixed cost inversion is not well-defined. Instead, I treat observed entry as data from a binomial process:
\begin{align}
    n_{r\tau t} \sim \mathrm{Binomial}(\bar{N}_{r\tau},p_{r\tau t}(\theta_\tau, \sigma_\tau)),
\end{align}
where $\theta_\tau$ are the non-monetary parameters of the fixed cost function and $\sigma_\tau$ is the scale of the cost shock that entrants receive. Then, I select $\theta_\tau$ such that model fixed costs maximize the binomial likelihood of the data:
\begin{align}
    \hat{\theta}_\tau(\sigma) = \argmax_{\theta_\tau} \sum_{r,t} [n_{r\tau t}\log p_{r\tau t} + (\bar{N}_{r\tau} - n_{r\tau t}) \log (1 - p_{r\tau t})].
\end{align}
I estimate this likelihood separately for solar and wind entrants. For solar entrants, the specification includes net capex shifters and an intercept $\alpha_{\mathrm{solar}}$, with no other only non-cost shifters. For wind entrants, I include an intercept $\alpha_{\mathrm{wind}}$ and a measure of incumbent wind stock $\log(1 + q_{r,\mathrm{wind}})$. Including a measure of existing stock for wind helps approximate the effect of omitted development difficulties like permitting issues. 

Both solar and wind thus produce $\hat{F}_{r\tau}$, fixed costs that can be represented by shifters and fixed effects. Comparing the baseline model results to counterfactuals requires using consistent fixed cost parameterizations across each scenario; to facilitate this, I re-estimate equilibrium using $\hat{F}_{r\tau}$. After re-estimating each fixed point, I am left with $\{p^\star, V^\star, M^\star(g|x),F^\star\}$.

Estimating $\sigma$ remains. To do so, I test a grid of different sigma values ranging from $\sigma=4$ to $\sigma=4096$. For each technology, fixed-cost specification, and candidate $\sigma$, I proceed as follows. 

I hold out one region $r$. Then, using all other regions, I estimate fixed cost coefficients:
\begin{align}
    \hat{\theta}_{\tau,-r}(\sigma) = \argmax_\theta \sum_{j \notin r,t} \mathcal{l}(n_{j\tau t}, \bar{N}_{j\tau}; p_{j\tau t}(\theta, \sigma)).
\end{align}
I then use the fitted fixed cost function to predict entry in $r$ and calculate the binomial log likelihood for the held out region. Then I repeat this until every region has been held out once and add together the likelihoods. I then select the $\sigma$ that maximizes the sum of the likelihoods.

%\subsection{Remark}

%It is worth discussing the estimation of the MME in this paper in the context of oblivious equilibrium and MME more broadly. Both oblivious equilibrium and MME evolved out of the need to avoid the curse of dimensionality in dynamic games. In an oblivious equilibrium, firms react to long-run expected industry states, and not realizations of that industry state. Firms are small and do not think about how their actions change industry states. \cite{IfrachWeintraub2017} introduce MME to deal with industries where large firms track each other's actions and a competitive fringe can discipline the large firms. My setting is closer to one where firms do not know their own position within the broader product market. 

\section{Results}
\label{sec:results}
This section details the results from the dynamic entry model. I describe the estimated fixed costs $\bar{F}_{r\tau}$ and the model-implied project values $V_{r\tau}$. I then look at the spatial allocation of entry for the model relative to the observed spatial allocation of entry. 

\subsection{Model Values}

Estimated project costs $\bar{F}_{r\tau}$ and model-implied project values $V_{r\tau}$ are shown in Table \ref{tab:dynamic_entry_estimates_values}. The table aggregates values across regions and across time periods. Costs and values are reported in MUSD/MW. I find that fixed costs for wind generators are almost three times larger than they are for solar generators. The distribution of fixed costs for solar is almost uniform across regions, whereas wind fixed costs vary significantly across regions. This is because of the fixed cost specification, where wind fixed costs are discounted in regions with pre-existing wind entry. Conceptually, this is consistent with a model where there are significant regional benefits to pre-existing wind, like developers already working in the area, streamlined permitting processes, and lesser interconnection fees. I assume that those regional benefits do not extend to solar generators.\footnote{I am testing a specification that includes pre-existing $q$ for solar as well, which I think will probably fit better and look more reasonable.} For both wind and solar, fixed costs fall during the period I study. For solar, this is driven by sharp capital cost decreases, while for wind it is primarily driven by tax benefits. 

Fixed cost results imply total fixed costs of \$111 MUSD for solar generators and \$421 MUSD for wind generators. A 2023 EIA report estimated installed construction costs for solar at \$196 MUSD and wind at \$294 MUSD respectively (\cite{EIA2025GeneratorConstructionCosts2023}).\footnote{\cite{SeelEtAl2024UtilityScaleSolar} and \cite{WiserEtAl2024LandBasedWind} estimate similar fixed costs for wind and solar.} Higher wind fixed costs are partially due to the structure of the fixed cost parameterization: wind fixed costs vary with pre-existing $q_{\mathrm{wind}}$ in a region and solar fixed costs do not, so regional fixed cost variation for wind evolves endogenously. Wind fixed costs are closer to estiamtes in the literature once I account for this effect -- fixed costs for wind generators in regions with $\geq 500$MW of wind generation are \$272 MUSD. 

\input{../figures/paper/entry_model_results/dynamic_entry_estimates_and_values.tex}

Table \ref{tab:dynamic_entry_estimates_values} also reports implied $V_{r\tau}$ for wind and solar. I find that mean $V_{r\tau}$ is lower than fixed costs for wind and solar, consistent with entry rates of 5\% and 13\% respectively. Mean $V_{r\tau}$ is roughly two-thirds the size of average fixed costs for solar and one-third the size of fixed costs for wind. Project values for wind are higher than they are for solar. Intuitively, $V_{r\tau}$ depends on (and increases in) the amount of electricity that a generator produces over its lifetime and the price it receives for that electricity. Wind generators have higher capacity factors (a measure how much electricity a generator can produce relative to its capacity) than solar generators do, which drives the observed differences in $V_{r\tau}$. This is true even though wind generators face lower prices than solar generators do: on average, LMPs for wind generators are 4.3\% below the system price of electricity, whereas they are only 0.9\% below the system price for solar generators.\footnote{All results in Table \ref{tab:dynamic_entry_estimates_values} are conditional on entry.} The discount faced by wind generators varies dramatically; in one region (Willacy county, in the far south of Texas), average congestion discounts faced by entrants are 24.7\%. 

Figure \ref{fig:entry_model_regional_fit} assesses the spatial fit of the dynamic entry model. Each point in the figure represented a model region. The x-axis shows the observed number of entrants for each region between 2018 and 2024, and the y-axis shows the model-implied number of entrants for that region between 2018 and 2024.\footnote{Including an intercept in the fixed effect parameterization mechanically recovers entry in the aggregate.} Overall fit is strong: correlation between model and observed entry for solar is 0.89 and for wind is 0.96. A line of best fit for each graph has slope 0.50 for solar, suggesting room for improvement, and 0.75 for wind, consistent a stronger overall fit. The model compresses the spatial distribution towards the mean for both wind and solar. 

\begin{figure}[!htbp]
    \centering
    \includegraphics[width=\textwidth]{../figures/paper/entry_model_results/entry_model_regional_fit_scatter.pdf}
    \caption{Observed and Model-Implied Cumulative Regional Entry, 2018--2024}
    \label{fig:entry_model_regional_fit}
    \fnote{\textit{Notes:} Each point is a model region. The dashed line is the 45-degree line.}
\end{figure}

There are several points where observed wind and solar entry is higher than model-predicted wind and solar entry. The far west is difficult for both wind and solar to pin down correctly: in the data, 18 wind generators and 23 solar generators entered the far west, while the model predicts only 13.5 and 11.5 generators enter in that region. One potential reason for this comes from Table \ref{tab:production-summary}: the far west is where demand grew most dramatically over this period, and thus where one might expect the stationary demand assumption to be more problematic. Other regional misses are less symptomatic: wind generator entry is too low in the eastern plains of Texas; solar entry is too low in west-central Texas.

To understand the magnitude of the congestion and resource quality channels, I estimate entry under two alternative scenarios. In the first scenario, I set congestion costs equal to zero everywhere on the grid and allow entry to adjust. Note that I do not estimate a new equilibrium. Figure \ref{fig:dynamic_entry_level_reallocation_decomposition} reports both the change in total entry and the change in its regional distribution. I find that wind entry increases by 8.9\%, a total of 8.2 project equivalents. The new distribution of entry is mostly centered on the original distribution of entry: the distance between baseline model entry shares and no-congestion entry shares is 1.9\%. Solar entry changes are more modest: total entry increases by 1.8\% (2.2 project equivalents) and the distance between shares is 1.1\%. Thus, eliminating congestion primarily affects the level of wind entry, with comparatively small changes in the spatial allocation of either technology.

Next, I examine how entry changes if resource quality were to shift. First, I give wind generators solar generators' resource quality distribution and geography, rescaled to preserve wind's maximum resource quality. I find this increases wind entry by 2.9\% (2.7 project equivalents), while shares differ by 2.8\%, indicating that adopting solar's resource geography substantively changes where wind generators locate. In the opposite exercise, adopting wind's resource geography reduces solar entry by 3.8\% (4.6 project equivalents) while shares differ by 7.0\%. Thus, wind's resource geography substantially changes where solar enters even though the decline in total solar entry is relatively modest.

\begin{figure}[!htbp]
    \centering
    \includegraphics[width=\textwidth]{../figures/paper/entry_model_results/dynamic_entry_level_reallocation_decomposition.pdf}
    \caption{Entry Levels and Spatial Distributions under Alternative Scenarios}
    \label{fig:dynamic_entry_level_reallocation_decomposition}
    \fnote{\textit{Notes:} Dark bars report partial-equilibrium changes in total entry relative to the estimated baseline. Light bars report the total-variation distance between the baseline and counterfactual regional entry-share distributions: one-half the sum of the absolute regional share changes.}
\end{figure}

\section{Counterfactuals}
\label{sec:counterfactuals}
I seek to answer two questions in counterfactuals: first, how do different infrastructure policies shift the equilibrium entry behavior of firms, and what are the consequences for welfare? And second, as demand for electricity grows, how do different infrastructure policies respond to increased demand, and what are the consequences for welfare? I test three different infrastructure policies: the ERCOT baseline, a slightly modified version of the PJM baseline, and an open-season style policy, where I set up an idealized market for new transmission.\footnote{Designing an incentive-compatible auction for new transmission is a promising direction for future research.} 

The modified version of the PJM baseline is designed to satisfy a reliability baseline: there must be sufficient transmission capacity for entering generators to fully deliver (i.e., deliver without congestion) the electricity they produce under three separate system conditions, which are designed to mimic periods of grid stress. If, for any of those system conditions, the entering generator would be unable to deliver its electricity, then the grid operator requires that generator to fund transmission improvements to ensure deliverability. 

The market for new transmission extends the PJM baseline in two directions. First, it allows entrants, incumbents, and demand to participate in the market. Second, it allows participants to submit bid curves: starting from their current level of deliverability in each planning case, the participant expresses their willingness to pay to increase that deliverability up to the PJM baseline (full deliverability).\footnote{Demand is assumed willing to pay the savings via lower prices resulting from the investment.} The grid operator then selects the upgrade package that maximizes willingness to pay minus upgrade costs. 

Transmission is fixed in the ERCOT baseline. The PJM baseline and the market for new transmission run each period. Upgrades are completed before dispatch and persist; entrants anticipate that the policy will run each year and make irreversible changes to the grid. 

\subsection{Model Additions}

The counterfactuals require two additions to the model: a method for selecting upgrades and state variables that contain information about grid capacity. I discuss each of them in turn. 

There are 300 corridors in the dispatch model, each of which is eligible for an upgrade. I treat upgrades as a flat 1200MW capacity addition to a corridor. There are significant complementarities across upgrades. For example, if a rural, high-wind region is connected to demand by a single transmission line that has two corridors, upgrading one of those corridors does not improve the grid. Congestion still arises on the unupgraded corridor. Instead, properly evaluating the value of an upgrade requires evaluating the effect of upgrading both links at the same time. 

To do so, I proceed with a beam search through upgrade packages (add cite). Beam search, which comes from computer science, is a graph and tree search algorithm that focuses on searching near high-value nodes. For example, consider a network with four corridors: $A, B, C, \text{ and } D$. A beam search begins by evaluating an upgrade on each individual corridor. It then identifies the $W$ top upgrades and moves on to the next level of evaluation. Consider for example, $W=2$, where $A$ and $C$ survive. Then the search evaluates upgrade pairings $AB, AC, AD, BC, CD$, but not $BD$. The search continues moving through levels until potential upgrade groups are exhausted.\footnote{Note that $BD$ is never evaluated in this example; exhausting potential upgrade pairings mean all pairings which include either $A$ or $C$. All three- and four-level groups including $A$ and $C$ would be evaluated.}

For the purposes of the counterfactual, only corridors with congestion or corridors which reduce congestion elsewhere following an upgrade are included in the beam search. Any corridor upgrade that only becomes valuable conditional on another upgrade occuring and that also is not originally congested will not be included in the search. Assessing a corridor in the beam search involves applying a 1200MW capacity addition to the corridor, re-computing the PTDF matrix, and seeing how electricity flows adjust to determine whether the upgrade changes congestion on the grid. The width $W$ of the search at each level is 100.

The upgrades chosen in equilibrium shift congestion and therefore prices. Including the grid as a state would be computationally infeasible. Instead, I model generators as caring about upgrades insofar as they change local prices, curtailment, and profits. This effectively treats the grid as a hidden state, similar to the way that entrants take other-region entry into consideration. As a result, counterfactual beliefs $M(g,Z|x,k)$ are formed over the spatial allocation of entry $g$ and the capacity of the grid $Z$. Beliefs now depend on state variables $x$ and $k$, where $k$ is [xxxx]. 

\subsection{Results}

\begin{itemize}
    \item PJM: Entrant observes upgrade cost, option to withdraw, leads to cascade
    \item Market: Only built if transmission line is funded 
    \item Market: Truthful bids, allocation of upgrade costs
    \item Welfare definition (including emissions)
\end{itemize}

\subsection{Data Center Experiment}

\section{Conclusion}
\label{sec:conclusion}

\clearpage
\bibliographystyle{aea}
%\bibliographystyle{plainnat}
\bibliography{bib_interconnect}

\end{document}
